\documentclass[referee,pdflatex,sn-basic]{sn-jnl}

\usepackage{graphicx}
\usepackage{booktabs}
\usepackage{amsmath}
\usepackage{amssymb}
\usepackage{array}
\usepackage{caption}
\usepackage{hyperref}
\usepackage{url}
\setcitestyle{numbers}

\title[An Audit of a Released Back-Translation Evaluator]{An Audit of a Released Back-Translation Evaluator under Donor-Replay Stress Tests: A Recipe-Conditioned Non-Reproduction Case Study}

\author*[1]{} % Authors retained for blind review
\affil*[1]{}

\begin{document}

\abstract{Back-translation (BT) is widely used to evaluate sign language production, yet a released evaluator checkpoint may yield decoded-text metric orderings that do not survive retraining, retrieval-pool controls, or reference replacement. On the 641-sequence released PHX-public pose subset, a constructed whole-donor replay stress test (TN-PURE-v1; not an actual SLRTP2025 leaderboard entry) scores 23.79 sacreBLEU under the released SLRTP2025 evaluator versus 12.78 for the recorded test poses, a gap of $+11.01$ BLEU points (95\% CI $[+9.61,+12.43]$). Fourteen recipe-conditioned reconstructions of the released pipeline do not reproduce this magnitude (gaps $-1.10$ to $+0.13$), and neither do train-pool-scaling variants, twenty rescue variants, all three examined checkpoint-selection objectives (NLL, decoded BLEU, decoded WER) over the training trajectories, or two step-faithful runs under the corrected released protocol. The reversal is highly sensitive to reference choice (attenuating to $+1.12$ against public human sign-to-text back-translations) and is confined to decoded BT-text metrics; teacher-forced BT likelihood does not reproduce it. Five related diagnostic views point to the same property: it decodes its full 7{,}060-item training pool at 78.4 BLEU with a 69.4\% exact-match rate while generalizing to unseen data (13.38/12.78), reproduces donor transcripts from donor motion at corpus BLEU 81.2 vs 23.8 (ratio 3.42), compared to 1.31--1.43 for reconstructions, loses $-16.1$ points when donor frames are reversed, shows a flat gap--competence dose--response within the observed competence range (dev BLEU-4 4.5--10.0, gap $[-1.1,+0.3]$), and its advantage grows with donor--query similarity; two frozen-protocol confirmation seeds reproduce the family side of this pattern. The evidence is consistent with the anomaly being a property of the released checkpoint and its (partly undocumented) data pipeline, co-occurring with unusually strong train-pool readout under the original reference frame; the specific cause is not identifiable from the released artifact. It is not a property of BT evaluation in general, not a leaderboard-rank statement, and not evidence that a randomly reseeded reconstruction would behave alike.}

\keywords{sign language production, back-translation evaluation, retrieval-augmented generation, multimedia quality assessment, evaluator non-replication}

\maketitle

\section{Introduction}

Sign language production (SLP) systems are commonly evaluated by back-translation (BT). A trained sign recognizer decodes each output pose sequence, and a text metric compares the decoding with the reference. This protocol measures recognizer-readable content through a learned measurement model. System rankings may therefore reflect both the evaluated outputs and the inductive biases of a particular evaluator checkpoint. When proposing BT evaluation for SLP, Saunders et al.~\cite{saunders2021continuous} argued that changing the BT model changes absolute scores but preserves the relative ranking of evaluated systems; the present audit directly tests a stress case of that rank-stability claim.

Previous studies show that BT scores can be modest for recorded real signing motion and can improve without faithful target conditioning \cite{walsh2024datadriven,hong2026conditional}. Sign-native and pose-aware metrics provide complementary evidence \cite{signbleu,silverscore,jiang2025meaningful}. Nevertheless, BT remains central to benchmarks such as SLRTP2025 \cite{slrtp2025}. An unresolved question is whether a retrieval--reference stress-test ordering persists across independently trained instances of the same BT architecture.

This question is exposed by a retrieval--reference reversal on PHX-public. Under the original SLRTP2025 evaluator, source-conditioned whole-donor retrieval scores 23.79 sacreBLEU (0--100 scale), whereas the recorded test poses score 12.78. This result neither establishes better signing by retrieval nor invalidates BT evaluation in general. Instead, it asks whether the reversal persists across evaluator checkpoints or depends on properties of the original evaluator.

We address this question in four stages. First, we compare the reversal with reference-pose, generative and random-donor controls under the original evaluator, verifying protocol details (scorer equivalence, frame-rate handling, decoding parameters) and aligning all secondary text metrics with the official repository implementations. Second, we evaluate the same fixed outputs with fourteen independently trained recipe-conditioned reconstructions and quantify paired uncertainty, documenting that no rescue variant, any of the three examined checkpoint-selection objectives (NLL, decoded BLEU, decoded WER), or two step-faithful runs under the corrected released protocol reaches competence parity. Third, we test reference validity directly: replacing the speech-transcribed PHOENIX references with public human sign-to-text back-translations~\cite{czehmann2026sacred} attenuates the reversal from $+11.01$ to $+1.12$ BLEU points, and a reference-anchored weather-slot diagnostic shows the replay's decoded content tracks the donor transcript. Fourth, we characterize what distinguishes the released pipeline: donor-pool substitution designs (path-dependent, so reported as descriptive contrasts alongside a duration/length/signer-matched 605-query factorial with residual similarity imbalance), and five related diagnostic views of a training-content readout property --- a training-pool readout signature combined with generalization that no reconstruction exhibits, donor-transcript pass-through far above the family trend, temporal-order sensitivity confined to the released evaluator, a flat gap--competence dose--response, and a frozen-protocol confirmation on two fresh seeds. Evaluator-aware oracle analyses are relegated to the appendix because their unmasked advantage does not survive reference-overlap exclusion. Together, these analyses distinguish a fixed-checkpoint, reference-dependent observation from a checkpoint-robust conclusion (Fig.~\ref{fig:audit_framework}).

\begin{figure}[t]
\centering
\includegraphics[width=\textwidth]{generated_figures/checkpoint_sensitivity_schematic.pdf}
\caption{Overview of the multi-checkpoint evaluation audit. (a) A fixed set of reference-pose, generative, retrieval, and random-donor outputs is evaluated by the original SLRTP2025 back-translation checkpoint and fourteen independently trained recipe-conditioned reconstructions. Because system outputs are held fixed, changes in scores, rankings, and retrieval--reference gaps reveal sensitivity to the evaluator instance while leaving competence and training-history differences entangled. (b) In the evaluator-aware analysis, donors selected by maximizing likelihood under one checkpoint are evaluated by all seven checkpoints, distinguishing within-checkpoint optimization from cross-checkpoint transfer.}
\label{fig:audit_framework}
\end{figure}


\section{Related Work}

\subsection{Learned evaluation in sign language production}
Neural SLP maps spoken language or glosses to pose and often evaluates the generated sequence with a frozen sign recognizer \cite{camgoz2018neural,saunders2020progressive}. BT was proposed as an SLP evaluation protocol by Saunders et al.~\cite{saunders2021continuous}, who reported that the choice of BT model shifts absolute scores while leaving relative system rankings unchanged in their experiments; our audit is a direct stress test of that rank-stability claim, using a measurement regime (whole-donor replay from the evaluator's training pool) outside the systems they compared. BT is useful because it measures recognizer-readable linguistic content, but a trained model mediates between pose and the final text metric. SignBLEU \cite{signbleu} operates on multichannel symbolic transcriptions, while SiLVERScore \cite{silverscore} and pose-native measures offer complementary views. BT BLEU nevertheless remains central to benchmarks such as SLRTP2025 \cite{slrtp2025}. Prior work has already questioned BT construct validity and the stability of learned diagnostics; our narrower contribution is a fixed-output replication of one retrieval--GT reversal across recipe-conditioned BT evaluator reconstructions, not the first study of learned-metric stability.

\subsection{Retrieval and shortcut risks}
Retrieval-based SLP ranges from phonetic or sign-unit composition to retrieval-enhanced generation \cite{walsh2024datadriven,signstitching,soke}; articulator-aware models further motivate separating body, hand, and face channels \cite{darslp}. The SLRTP2025 report also describes a leading retrieval-based entry, making donor copying pertinent to the benchmark's evaluation paradigm \cite{slrtp2025}. More generally, memorization and retrieval-augmented language modeling show that strong task scores can reflect reproduction of training examples rather than synthesis \cite{carlini2021extracting,kandpal2022memorization,khandelwal2020knnlm}. In SLP, inserting real donor motion creates the analogous risk that a recognizer rewards lexical overlap with the donor. Donor exclusion and random-donor controls can localize that risk, but a shortcut observed under one learned evaluator does not show that the resulting ordering persists under another checkpoint.

\subsection{Robustness of learned metrics and corpus audits}
Existing analyses expose several kinds of metric fragility. Hong and Ko\v{s}eck\'{a}~\cite{hong2026conditional} (recent preprint) show that BT-BLEU and motion metrics can diverge from target faithfulness; Walsh et al.~\cite{walsh2024datadriven} document a low ground-truth BT baseline; and Jiang et al.~\cite{jiang2025meaningful} compare pose-based metrics with human judgment and recommend BT likelihood over decoded BLEU. Yazdani et al.~\cite{yazdani2025critical} test text metrics under paraphrase, hallucination, and length perturbations, while BackTranslation2.0~\cite{backtranslation2} proposes a linguistically grounded alternative to na\"ive BT evaluation (recent preprint). He et al.~\cite{he2023blind} expose systematic blind spots of model-based text-generation metrics under distribution shift, motivating audit-style stress tests of learned scorers. Semantic slot-level checks (SLU-2K~\cite{slu2k}) provide a complementary factual-content criterion, which we adapt as a rule-based weather-slot extractor (\S\ref{sec:slot_f1}). Two 2026 PHOENIX-2014T audits further constrain interpretation: Czehmann et al.~\cite{czehmann2026sacred} released human sign-to-text back-translations of the test set and documented information loss, lexical errors and translationese in the original speech-transcribed references; Alkain et al.~\cite{alkain2026critical} showed PHOENIX-2014T has the highest train--test text overlap among tested SLT corpora (about 5\% of test sentences appear verbatim in training) and that corpus BLEU rises steeply with training-likeness. Artiaga et al.~\cite{artiaga2025signer} showed signer-dependent evaluation overestimates SLT capability (BLEU-4 dropping from 21.44 to 3.59 for GFSLT-VLP and from 22.74 to 3.66 for SignCL under signer-fold cross-validation). None of these studies runs multi-checkpoint evaluator experiments or donor-pool substitution controls; our contribution is a fixed-output, multi-checkpoint audit of one released evaluator pipeline, connecting their corpus-level findings to the evaluator-level anomaly.

\section{Methods}
\label{sec:method}

\subsection{Evaluation design}
Let $E_c$ denote a BT evaluator checkpoint and $o_{i,n}$ the pose sequence for output condition $i$ and evaluation item $n$. Checkpoint $c$ produces
\[
h_{c,i,n}=E_c(o_{i,n}), \qquad
B_{c,i}=M\!\left(\{(h_{c,i,n},y_n)\}_{n=1}^{N}\right),
\]
where $y_n$ is the reference text and $M$ is the corpus-level sacreBLEU functional unless stated otherwise. The values $B_{c,i}$ induce a ranking over fixed output conditions. We call a result \emph{evaluator-pipeline-sensitive} when a substantive pairwise ordering changes across independently trained checkpoints under a fixed output set and protocol. The fixed primary set comprises the recorded reference poses (GT-pose, or GT for short), random donor copy, source-conditioned whole-donor retrieval, a Progressive Transformer generator, and retrieval--generator compositions; oracle-gloss variants are diagnostic controls.

The main evaluation uses PHX-public, a German weather-domain text-to-pose benchmark. The official metadata contain 7,096/519/642 train/development/test IDs. The released SLRTP2025 pose records available to this study contain 7,060/515/641 sequences: 36/4/1 official IDs are absent from the released pose materialization. Their IDs and the reason \texttt{absent\_from\_slrtp2025\_released\_pose\_records} are enumerated in \texttt{artifact/manifests/exclusions.jsonl}. The missing test ID is \texttt{13December\_2010\_Monday\_tagesschau-6940}; its text/gloss metadata are not present in the released pose record, so length and template density cannot be reconstructed without an additional official annotation source. We therefore condition every result on the complete 641-record release rather than assert that the missing item is ignorable. As an empirical observed-extreme sensitivity scenario, assigning the missing item the per-item score of the largest PURE-over-GT outlier observed among the 641 released items changes the corpus-level gap by less than $\pm 0.2$ BLEU points, well within the reported 95\% paired-bootstrap interval $[+9.6,+12.4]$. This is an observed-extreme scenario, not a worst-case bound: corpus BLEU is non-linear and the missing item's reference, hypothesis, and n-gram sufficient statistics are unknown, so the true mathematical worst case is not covered. All primary conclusions are therefore restricted to the released 641-record materialization. The official SLRTP2025 leaderboard reports the same GT BLEU-4 of 12.78 on the public test subset, consistent with the 641-item evaluation at 12.78 BLEU points. The supplementary analyses use two template-disjoint controls with distinct layouts and sizes: \textbf{PHX-holdout-178}, a 178-joint split with 525 held-out sequences used in the checkpoint experiments; and \textbf{PHX-holdout-MSKA}, a frozen MSKA split with 528 held-out sequences (1,538 windows) in a 133-joint layout, whose separate 499-sequence development subset is used only to check Total Distance. CSL-Daily provides a cross-dataset boundary control with 58{,}921 training windows and 256 test windows. The anonymous supplementary artifact contains scorer implementations, manifests, decoded hypotheses, sufficient statistics and secondary protocol records cited below. All paths reported in the paper are relative to its root.

\begin{table}[t]
\centering
\scriptsize
\caption{Experimental protocol summary. PHX-public denotes the SLRTP2025 public test, PHX-holdout-MSKA the frozen 133-joint template-holdout split, and CSL the CSL-Daily 256-window protocol.}
\label{tab:protocol_summary}
\begin{tabular}{llllrr}
\toprule
Setting & Recognizer & Layout & Split size & $N$ eval & $\alpha$ \\
\midrule
PHX-public  & SLRTP2025 evaluator       & 178-joint & 641 seq      & 641    & 0.88 (tuned) \\
PHX-holdout-MSKA & mska\_slr (phoenix-2014t) & 133-joint & 1{,}538 win  & 1{,}538 & 0.50 (frozen) \\
CSL-Daily   & mska (csl-daily)           & 133-joint & 256 win      & 256    & 0.50 (frozen) \\
\bottomrule
\end{tabular}
\end{table}

All pose sequences are evaluated at 12.5\,fps, matching the SLRTP2025 evaluator training distribution (\texttt{skeleton\_subsample=2}). Two protocol facts shape every experiment in this study. First, every canonical pose bundle (GT, retrieval, composed) is \emph{stored} at 12.5\,fps after exactly one construction-time \texttt{[::2]} transform (Table~\ref{tab:protocol_regression}), and the audit bypasses the released wrapper (which would subsample 25-fps inputs again) and calls the BT model directly, so no double subsampling occurs; a three-path regression confirms it, since the canonical stored-12.5-fps path and an explicit 25$\to$12.5-fps path give identical GT BLEU-4 (12.78) and official WER (85.77), while an erroneous second subsample (6.25\,fps) degrades them (6.65 / 86.85). Second, the released BT model applies \emph{no} internal subsampling: its \texttt{subsample} attribute is set but never used in the released code, so any evaluation path that feeds raw 25-fps poses directly is off-protocol by a factor of two. Every input in this study is therefore subsampled externally, exactly once.

\begin{table}[h]
\centering
\scriptsize
\caption{Protocol regression (left: all systems are stored and evaluated at 12.5\,fps after exactly one construction-time \texttt{[::2]}) and three-path FPS sensitivity (right: GT under the original evaluator; WER official, 0--100). Paths A and B agree exactly; the double-subsample error path C degrades both metrics, so the canonical result is not an artifact of double subsampling.}
\label{tab:protocol_regression}
\begin{tabular}{llrr}
\toprule
System & Source material & Stored fps & Stored frames \\
\midrule
GT (released test.pt) & 25 fps, 64{,}504 fr. & 12.5 & 32{,}411 \\
PT-v1 (Codabench starter) & 12.5 fps bundle & 12.5 & 32{,}400 \\
TN-PURE-v1 / TN-PTCOMP-v1 & 25 fps donors & 12.5 & 30{,}711 / 32{,}400 \\
SEEN-MATCHED / RAND640 & 25 fps donors & 12.5 & 29{,}418 / 32{,}708 \\
UNSEEN-PURE / CROSSFIT-A/B & 25 fps donors & 12.5 & 30{,}027 / 15{,}909 / 14{,}187 \\
\midrule
\multicolumn{4}{l}{\textit{Three-path (GT): A canonical 12.78/85.77; B explicit 25$\to$12.5 12.78/85.77;}} \\
\multicolumn{4}{l}{\textit{C double-subsample (error) 6.65/86.85}} \\
\bottomrule
\end{tabular}
\end{table}

The official Codabench starter-bundle Progressive Transformer (\texttt{PT\_baseline\_test.pt}) scores 1.59 BLEU points under the released evaluator on byte-identical tensors, matching the official 1.5942. (The challenge paper's 5.43 is a separately trained challenge-entry PT, not the starter-bundle baseline we audit as the weak-generation control.)

PHX-public uses $\alpha{=}0.88$, selected on its development set by targeting Total Distance near 1.0; PHX-holdout-MSKA and CSL-Daily use the pre-specified $\alpha{=}0.5$. Pose normalization, layout compatibility, and scaffold implementation are documented in the anonymous review artifact.

\subsection{Systems and retrieval controls}
\label{sec:retrieval}
For each target, source-conditioned whole-donor retrieval scans the complete training pool after Unicode NFKC normalization, whitespace normalization and lowercasing, and selects the sequence with maximum source-text token-set Jaccard similarity. Oracle-gloss retrieval instead uses target gloss and is treated as an input-leakage diagnostic; random retrieval selects a donor uniformly. Donor exclusion removes exact normalized text or gloss duplicates or donors above threshold $\tau$. Gloss-based exclusion uses target gloss and is therefore an offline diagnostic rather than a deployable inference procedure.

\paragraph{Text-nearest tie-breaking.}
\label{sec:method_tie_breaking}
Jaccard ties are resolved by minimum character-level Levenshtein distance and then by minimum SHA-256 hash of the training key. This deterministic rule is used throughout; alternative tie-breaking rules are included in the reproducibility material.

\paragraph{Retrieval composition.}
\label{sec:composition}
The canonical whole-donor output TN-PURE-v1 downsamples the selected 25-fps donor by taking every second frame and otherwise preserves its full duration. The distinct canonical composition TN-PTCOMP-v1 uses the same frozen donor registry, downsamples the donor to 12.5 fps, and linearly resamples it to the PT scaffold duration. The composed output combines a generator scaffold $B$ with a donor $R$. Hips and the remaining body follow $B$, whereas hands and face are copied from $R$. Upper-body joints use $Y[t,\mathcal{U}]=R[0,\mathcal{U}]+\alpha(R[t,\mathcal{U}]-R[0,\mathcal{U}])$. For SLRTP178, $\mathcal{U}{=}\{0,\ldots,5\}$, $\mathcal{A}{=}\{8,\ldots,177\}$ and $\mathcal{O}{=}\{6,7\}$. The corresponding MSKA sets are $\mathcal{U}{=}\{5,\ldots,10\}$, $\mathcal{A}{=}\{23,\ldots,132\}$ and $\mathcal{O}{=}\{11,12\}$. PHX-public uses Progressive Transformer output as its scaffold; layout-incompatible settings use the supplementary diffusion scaffold documented in the artifact. Whole-donor retrieval uses neither a scaffold nor $\alpha$. For the canonical registry, query and selected-donor durations match closely (median 3.96 s vs 3.76 s; ratio 0.957) and word counts match (median difference 0); full distributions are in the artifact. In the separately matched seen--unseen factorial, every retained four-cell set is matched on both duration and word count; 605/641 targets satisfy that stricter design.

\paragraph{Donor-origin contrast retrieval systems.}
\label{sec:exposure_systems}
To test whether the released-evaluator retrieval advantage co-occurs with donor origin, we construct frozen probe systems from the same 641 queries: \textbf{UNSEEN-PURE-v1} retrieves from the other PHX-public test poses (640 candidates after excluding the query itself, its sequence root, and normalized exact-text matches); \textbf{CROSSFIT-PURE-A/B-v1} split the queries by hash into folds of 331/310, each retrieving from the other fold's poses; \textbf{UNSEEN-MATCHED-v1} is the 302-query subset of UNSEEN-PURE-v1 admitting a SEEN donor within Jaccard tolerance 0.1 and duration tolerance 1\,s; \textbf{SEEN-PURE-RAND640-v1} draws a fresh 640-donor random train sub-pool per query (pool-size control); \textbf{SEEN-PURE-MATCHED-v1} selects, for each query with an UNSEEN donor, a train-pool donor within $\pm0.1$ of the UNSEEN donor's Jaccard (622/641 admit one; mean $|\Delta J|=0.006$); and \textbf{SEEN-RAND640-MATCHED-v1} applies the matched selection inside the frozen 640-donor sub-pools (mean $|\Delta J|=0.052$; 15.1\% of queries have $|\Delta J|>0.1$ because the 640-donor sub-pool limits match quality). Signer is not constrained in the primary matched design (the \texttt{speaker} field is populated for all released records); a signer-strict variant is reported in \S\ref{sec:balance}. SEEN-PURE-v1 re-uses the canonical TN-PURE-v1 registry. These systems are constructed probes, not deployable generators: UNSEEN-PURE-v1 retrieves from the test set itself.

\subsection{Evaluator checkpoint family}
The primary inferential checkpoint family comprises the original frozen SLRTP2025 evaluator and six \emph{recipe-conditioned reconstruction} runs trained from scratch with the documented architecture and a common frozen train/development partition. Primary seeds are 101, 202, 303, 404, 505 and 606; eight extension seeds (707, 808, 909, 1001, 1102, 1203, 1304 and 1405) use the identical protocol; only initialization and seeded data order vary. Checkpoints are selected by the same validation-loss rule. We call these runs \emph{recipe-conditioned reconstructions} (formally, the \emph{paper-derived reconstruction protocol} of Table~\ref{tab:recipe_provenance}) rather than replicas or IID seed repeats. The original bundle contains more than the checkpoint: \texttt{best.ckpt}, \texttt{config.yaml}, the two vocabularies, and \texttt{validations.txt} --- and we parse that log in \S\ref{sec:original_trajectory}: it records 202 validation points over 2{,}828 optimizer steps (validation every 14 steps; at batch 256 over 7{,}060 training sequences this is $\approx$102 epochs), with decoded-BLEU selection and a best dev BLEU-4 of 13.38 at step 1{,}820 ($\approx$epoch 66), matching our independently re-measured dev BLEU-4 of 13.38 to two decimals. The released checkpoint is therefore the step-1{,}820 weights of a $\sim$100-epoch decoded-BLEU-selected run --- not a multi-thousand-epoch run --- while intermediate checkpoints are not published. Every checkpoint evaluates byte-identical canonical pose outputs, identified by semantic tensor SHA-256. The full reconstruction recipe (architecture, tokenizer, vocabularies, train/dev manifests, optimizer, scheduler, label smoothing, dropout, batching, checkpoint-selection rule, dependency versions, and training commit) is enumerated in the artifact's protocol file and machine-readable contract. Table~\ref{tab:recipe_provenance} maps each reconstruction setting to its source; we call the fourteen main runs the \emph{paper-derived reconstruction protocol} (the training loop follows the documented legacy script: 300 epochs, validation-NLL selection) and the four runs of \S\ref{sec:config_faithful_results} the \emph{released-config reconstruction protocol} (following the released \texttt{config.yaml}: up to 3{,}000 epochs, validation every 14 epochs, decoded-BLEU early stopping). The 300-epoch/BLEU figure reported in the challenge paper refers to the Progressive Transformer production baseline, not to the BT evaluator, so the BT evaluator's actual training recipe remains unpublished beyond the released \texttt{config.yaml} and code.

\begin{table}[h]
\centering
\scriptsize\setlength{\tabcolsep}{3pt}
\caption{Reconstruction recipe provenance. Sources: P = challenge paper, C = released \texttt{config.yaml}, S = released source code, I = author inference from released artifacts, A = author-chosen for the reconstruction.}
\label{tab:recipe_provenance}
\begin{tabular}{llll}
\toprule
Setting & paper-derived protocol (14 runs) & released-config protocol (4 runs) & Source \\
\midrule
Architecture (3L/256/512ff, 8 heads) & same & same & C, S \\
Tokenizer / vocabularies & same & same & C, S \\
Train/dev manifests (7{,}060/515) & same & same & C (data paths), released .pt \\
Input subsampling ($[::2]$) & same & same & C (skeleton\_subsample=2), S \\
Optimizer Adam 1e-3, betas .9/.998, wd .001 & same & same & C, S \\
Batch size 256, grad clip 1.0 & same & same & C, S \\
LR schedule (plateau $\times0.8$) & same & same & C (scheduling, decrease\_factor), S \\
Epoch budget & 300 & up to 3{,}000 & A (bounded) / C (epochs) \\
Validation frequency & every 5 epochs & every 14 epochs & S (legacy loop) / C (validation\_freq) \\
Selection objective & dev NLL & decoded dev BLEU & I (legacy script) / C (early\_stopping\_metric=bleu) \\
Early-stop patience & 15 validations & 5 validations & A (bounded) / C (patience) \\
Seeds & 101--606, 707--1405 & 101--404 & A \\
\bottomrule
\end{tabular}
\end{table}

\subsection{Metrics and uncertainty}
The recorded test poses (GT) are decoded by each BT recognizer and compared with the original reference text. We report corpus BLEU-4 on the conventional 0--100 sacreBLEU scale (BLEU points) as the primary metric (e.g., 23.79 BLEU points) \cite{post2018sacrebleu}; sentence-level statistics (pass-through, slot-F1) are reported separately on the 0--1 scale and always labelled as sentence-level means. BLEU-1, chrF (0--100 scale, official corpus-level implementation), ROUGE-L (official script, 0--100), BERTScore \cite{zhang2020bertscore} (\texttt{bert-score} 0.3.12, \texttt{bert-base-multilingual-cased}, default layer, \texttt{lang=de}, \texttt{idf=False}, no baseline rescaling, F1), BLEURT \cite{sellam2020bleurt} (\texttt{Elron/bleurt-large-512}, a PyTorch-native substitute for BLEURT-20; Pearson $r=0.87$, Spearman $\rho=0.81$ on 100 German (ref, hyp) pairs against canonical BLEURT-20 drawn uniformly at random with a fixed seed from the PHX-public decoded cells before any comparison was run, with bias $-0.47$ so absolute scores are not interchangeable but rank-based conclusions are stable), and WER (official jiwer 3.1.0 normalized implementation, 0--100 scale) provide complementary decoded-text controls; BT likelihood and pose-DTWp provide non-decoded controls. The canonical BLEURT-20 checkpoint is distributed only as a TensorFlow saved\_model whose internal XLA JIT compilation fails under our TF-GPU setup; full consistency-validation numbers are reported in the artifact.

\paragraph{Scorer equivalence with the official SLRTP2025 scorer.}
The vendored official scorer (SacreBLEU 1.4.2 \texttt{raw\_corpus\_bleu}: \texttt{tok:none}, \texttt{smooth:floor}, \texttt{smooth\_value=0}, \texttt{use\_effective\_order=True}, \texttt{force=True}) and the paper scorer (SacreBLEU 2.5.1, \texttt{tok:13a}, \texttt{smooth:exp}, \texttt{effective\_order=False}, \texttt{force=True}) give \emph{identical} corpus BLEU on all 28 evaluation cells (7 evaluators $\times$ 4 systems; maximum absolute difference $0.0000$, and the same PURE--GT ordering under all 7 evaluators): the BT-decoded hypotheses are pre-tokenized and all four n-gram orders have non-zero precision, so neither tokenization nor smoothing differentiates. We report the 2.5.1/13a signature as primary and provide both outputs in the artifact. chrF/WER/ROUGE-L are reported under the official implementations (\S\ref{sec:metric_protocol_alignment}).
\begin{table}[h]
\centering
\small
\caption{Scorer equivalence on byte-identical decoded hypotheses: official SLRTP2025 scorer (SacreBLEU 1.4.2 \texttt{raw\_corpus\_bleu}) vs paper scorer (2.5.1, \texttt{tok:13a}/\texttt{smooth:exp}/\texttt{eff:no}). Maximum BLEU-4 difference across all 28 cells is $0.0000$; the PURE--GT ordering is identical under every evaluator.}
\label{tab:scorer_equivalence}
\begin{tabular}{lrrr}
\toprule
Quantity & Official (1.4.2) & Paper (2.5.1) & $|\Delta|$ \\
\midrule
Original GT BLEU-4        & 12.78 & 12.78 & 0.00 \\
Original TN-PURE BLEU-4   & 23.79 & 23.79 & 0.00 \\
Max $|\Delta|$ across 28 cells & --- & --- & $\mathbf{0.0000}$ \\
\bottomrule
\end{tabular}
\end{table}

We report two forms of uncertainty. \emph{Item-sampling uncertainty} uses paired sequence-level bootstrap (10{,}000 resamples, seed 42, percentile 95\% CIs, including the common-support decomposition on the full 622-query support): per item, n-gram sufficient statistics are precomputed once, and every resample re-aggregates them into a corpus BLEU (never an average of per-item sentence BLEUs). A pairwise difference is called detectable when its interval excludes zero. Significance testing of the primary estimand uses a paired approximate-randomization permutation (two-sided; the exchange unit is the per-item pair of decoded hypotheses, swapped with probability 1/2). Because the headline gap was first observed in exploratory analysis and the estimand was frozen afterwards, we report the result as a descriptive Monte-Carlo estimate rather than a confirmatory $p$-value: 0 of 10{,}000 permutations exceed the observed gap, i.e., $p\leq(0{+}1)/(10{,}000{+}1)\approx10^{-4}$ with an exact 95\% binomial interval $[0,0.00037]$ for the underlying rate (the lower bound is 0 by construction of the estimate). For the six primary runs we report difference-in-differences effect sizes with paired-bootstrap CIs but no bootstrap $p$-values (a paired bootstrap approximates the sampling distribution, not the null). \emph{Training-replicate uncertainty} is reported only for the recipe-conditioned population: a descriptive seed-level 95\% Student-$t$ prediction interval, $\bar{x}\pm t_{0.975,n-1}\,s\sqrt{1+1/n}$ (not extrapolable to the released evaluator, which is not a draw from that population). The eight extension seeds are post-hoc additions and enter only that interval and the descriptive ranges, never the six primary comparisons. The single primary estimand is the PURE--GT decoded-sacreBLEU gap under the released evaluator; all other contrasts are secondary and are Holm-corrected within a defined family or explicitly labelled descriptive (Appendix~\ref{sec:provenance}).

\section{Results}
\subsection{Retrieval reverses the reference ranking under the original evaluator}
\label{sec:public-diagnostic}
Under the original SLRTP2025 evaluator, source-conditioned whole-donor retrieval scored 23.79 sacreBLEU (0--100 scale), compared with 12.78 for the recorded test poses. The retrieval--ground-truth gap was $+11.01$ BLEU points (95\% paired-bootstrap CI, $[+9.61,+12.43]$ BLEU points; 10,000 sequence-level resamples). This fixed-checkpoint reversal motivates the analysis but does not indicate superior sign production.

Controls localized the observation to source-conditioned donor retrieval. The pure donor-copy row exceeded the oracle-gloss composed variant (11.8 BLEU points), the canonical text-nearest PT-composed variant (18.86 BLEU points), the Progressive Transformer baseline (1.59 BLEU points), and the random-donor baselines (0.90 BLEU points PURE and 0.77 BLEU points COMP, each over 20 common seeds; Table~\ref{tab:isomorphic_conditioning}). Source text was a legal inference input, whereas oracle-gloss retrieval used reference gloss and did not exceed the recorded test poses (11.4--11.8 vs 12.78 BLEU points). Token-normalized teacher-forced NLL (negative log likelihood per token; lower is better) does not reproduce the decoded-BLEU reversal under Original: GT is 3.060, compared with 3.479 for PURE, 3.411 for COMP and 4.074 for PT. The compatible normalized pose-DTW diagnostic likewise does not validate semantics: GT is 0 by construction, while PT, text-nearest PURE and random donor score 0.00779, 0.00830 and 0.00678, respectively (lower is better). The reversal is also visible under complementary text metrics computed on the same decoded hypotheses (Table~\ref{tab:metric_sensitivity}), here reported under the \emph{official} SLRTP2025 metric implementations after the protocol alignment of \S\ref{sec:metric_protocol_alignment}: under the original evaluator, sacreBLEU-4 reverses by $+11.01$ BLEU points, BLEU-1 by $+21.03$ BLEU points, official chrF by $+13.23$ points, official ROUGE-L by $+17.87$ points, official normalized WER by $-13.90$ points (negative is better), BERTScore-F1 (multilingual BERT) by $+0.066$, and BLEURT (\texttt{Elron/bleurt-large-512}) by $+0.165$. The decoded-text metrics agree on direction and rank; teacher-forced likelihood and geometric pose proximity do not. We therefore describe the observation as a \emph{decoded-text-metric reversal} rather than a BT-likelihood or pose-geometry reversal, and we use sacreBLEU-4 as the primary decoded metric throughout. Full per-checkpoint likelihood, pose, and complementary text metrics are supplied in the artifact.

\paragraph{Metric implementations.}
\label{sec:metric_protocol_alignment}
All secondary metrics use the official SLRTP2025 implementations: on byte-identical decoded hypotheses they reproduce the official repository values for the recorded test poses (BLEU-4 12.777, chrF 34.585, WER 85.775, ROUGE 35.196 \cite{slrtp2025}) \emph{exactly}. The implementation details matter materially for the two non-BLEU metrics, so we document the signatures: official chrF is corpus-level sacrebleu \texttt{CHRF} defaults (character n-grams 1--6, whitespace excluded), and plausible alternatives differ substantially --- a sentence-level variant over character n-grams 1--4 including whitespace yields 48.05 instead of 34.59 on the same hypotheses (full corpus/sentence $\times$ n-gram-order $\times$ whitespace ablation in the artifact). Official WER uses jiwer 3.1.0 with the official normalization (lowercasing, punctuation removal); a raw whitespace-split WER yields 79.26 instead of 85.78, because the normalization removes near-deterministic trailing-period tokens from the reference-word denominator (a raw-split WER on punctuation-stripped lowercase tokens reproduces 85.775 exactly). BLEU-4 is identical under all these variants, but chrF and WER are not, so we fix the official implementations throughout; alternative-protocol values are retained in the artifact as a sensitivity check.

\begin{table}[t]
\centering
\scriptsize\setlength{\tabcolsep}{3pt}
\caption{Decoded-text metric sensitivity under the original SLRTP2025 evaluator (641 PHX-public sequences; canonical decoding: beam size 3, length penalty $-1$, max output length 30). All metrics are computed on the same decoded hypothesis--reference pairs under the \emph{official} SLRTP2025 implementations (\S\ref{sec:metric_protocol_alignment}): chrF (sacrebleu \texttt{CHRF} defaults, corpus-level), ROUGE-L (official script), WER (jiwer 3.1.0 with the official normalization, 0--100 scale). BERTScore uses \texttt{bert-base-multilingual-cased} (bert-score 0.3.12, \texttt{lang=de}, \texttt{idf=False}, no baseline rescaling); a BLEURT substitute column is reported in the artifact as rank-based sensitivity. The reversal direction is consistent across decoded-text metrics; teacher-forced NLL and pose-DTWp (text) do not reproduce it.}
\label{tab:metric_sensitivity}
\begin{tabular}{lrrrrrrr}
\toprule
System & BLEU-4 & BLEU-1 & chrF & ROUGE-L & WER & BERTSc. \\
\midrule
GT-v1 (reference poses) & 12.78 & 34.40 & 34.59 & 35.20 & 85.77 & 0.757 \\
PT-v1         & 1.59 & 16.91 & 19.03 & 17.26 & 101.68 & 0.669 \\
TN-PTCOMP  & 18.86 & 46.35 & 41.88 & 45.17 & 77.36 & 0.793 \\
TN-PURE    & \textbf{23.79} & \textbf{55.43} & \textbf{47.82} & \textbf{53.07} & \textbf{71.87} & \textbf{0.822} \\
\midrule
Gap PURE$-$GT & $+11.01$ & $+21.03$ & $+13.23$ & $+17.87$ & $-13.90$ & $+0.066$ \\
\bottomrule
\end{tabular}
\end{table}

\paragraph{Decoding-parameter sensitivity.}
Re-decoding all four canonical systems under the original evaluator and seed 101 across beam size $\{1,3,5\}$, length penalty $\{-1,0\}$, and max output length $\{30,50\}$ shows the reversal is not an artifact of decoding configuration: under the original evaluator the PURE--GT gap is $+9.8$ (greedy), $+11.0$ (canonical beam 3), and $+11.6$ (beam 5), with length penalty and max length negligible ($<0.1$), while seed 101 stays in $[-0.4,+0.4]$ throughout (figure and full grid in the artifact).

\begin{table}[t]
\centering
\small
\caption{PHX-public conditioning and construction controls (641 sequences; 12.5\,fps). Columns distinguish whole-donor copy from PT-scaffold composition ($\alpha{=}0.88$); all rows share the query set, donor pool, evaluator, and sacreBLEU protocol. Semantic retrievers (LaBSE 19.2/8.4, multi-SBERT 16.5/8.1) lie between the canonical retriever and oracle gloss; full rows in the artifact.}
\label{tab:isomorphic_conditioning}
\begin{tabular}{lrr}
\toprule
Conditioning & Pure donor copy & PT+retrieval composed \\
\midrule
Source text (German Jaccard; canonical) & \textbf{23.79} & 18.86 \\
Oracle gloss (input leakage) & 11.4 & 11.8 \\
Random donor (seeds 0--19) & 0.90 (SD 0.24) & 0.77 (SD 0.19) \\
\bottomrule
\end{tabular}
\end{table}

Pure donor copy outperformed the composed variant for every source-conditioned retriever, and source-text Jaccard produced the highest pure-donor score (23.79, ahead of LaBSE 19.2 and multi-SBERT 16.5), while oracle-gloss retrieval did not exceed the recorded test poses in either construction (11.4--11.8 vs 12.78): the reversal does not require oracle-gloss input leakage, and pure copy --- not composition --- is its strongest form.



\subsection{Independent evaluator reconstructions do not reproduce the reversal magnitude}
\label{sec:what-drives-reversal}
TN-PURE-v1 yields an original-evaluator retrieval--reference gap of $+11.01$ BLEU points (95\% paired-bootstrap CI $[+9.61,+12.43]$ BLEU points). Across fourteen reconstruction runs, gaps range from $-1.10$ to $+0.13$ BLEU points; twelve point estimates are negative and two are small positive values near zero (seeds 101 and 606). TN-PTCOMP-v1 yields an original gap of $+6.09$ BLEU points ($[+4.76,+7.41]$ BLEU points); all fourteen reconstruction point estimates are negative. Thus no reconstruction reproduces the original effect magnitude. The seed-level 95\% prediction interval across the fourteen reconstruction Best checkpoints is $[-1.4,+0.3]$ BLEU points (PURE gap mean $-0.55$ BLEU points, SD $0.39$ BLEU points, 2 of 14 positive); this captures training-replicate variability that the sequence bootstrap cannot. The single released Original remains a case study rather than a draw supporting population inference. For the six primary runs, original-minus-run difference-in-differences estimates range from $+10.89$ to $+11.99$ BLEU points; every 95\% interval excludes zero. We do not report bootstrap $p$-values for these contrasts because a paired bootstrap approximates the sampling distribution of the statistic, not the null distribution.

\subsubsection{Recipe-conditioned evaluator reconstructions}
\label{sec:multiseed_factorial}
All fourteen reconstruction runs use the same frozen train/development manifests, architecture, hyperparameters and validation-loss checkpoint-selection rule; seeds 101--1405 vary initialization and data order. Table~\ref{tab:multiseed_factorial} reports the two canonical retrieval constructions separately.

\begin{table}[h]
\centering
\scriptsize
\caption{Canonical fixed-output evaluation on the released 641-pose PHX-public subset. PURE is TN-PURE-v1; COMP is TN-PTCOMP-v1. Gaps are relative to local GT (BLEU points). The six primary seeds carry precomputed paired-bootstrap intervals; the eight extension seeds (bottom row, ranges) reproduce the same near-zero or negative pattern post hoc. Per-cell values in the artifact.}
\label{tab:multiseed_factorial}
\begin{tabular}{lrrrrr}
\toprule
Evaluator & GT & PURE & Gap & COMP & Gap \\
\midrule
Original & 12.78 & 23.79 & $\mathbf{+11.01}$ & 18.86 & $+6.09$ \\
Seed 101 & 8.57 & 8.65 & $+0.08$ & 7.92 & $-0.65$ \\
Seed 202 & 8.81 & 7.83 & $-0.98$ & 7.54 & $-1.28$ \\
Seed 303 & 9.67 & 8.96 & $-0.71$ & 8.71 & $-0.96$ \\
Seed 404 & 7.19 & 6.83 & $-0.36$ & 6.38 & $-0.81$ \\
Seed 505 & 9.05 & 8.16 & $-0.89$ & 7.97 & $-1.08$ \\
Seed 606 & 8.09 & 8.22 & $+0.13$ & 7.93 & $-0.16$ \\
\midrule
Extension 707--1405 (range) & 7.71--9.56 & 6.84--9.09 & $-1.10$ to $-0.19$ & 6.90--8.83 & $-1.03$ to $-0.44$ \\
\bottomrule
\end{tabular}
\end{table}

Competence is assessed on the released 515-sequence development subset, never on PHX-public for checkpoint selection. Original obtains dev BLEU-4 13.38 and WER 83.37 (official jiwer implementation, 0--100 scale); reconstruction runs obtain 6.69--9.07 and WER 88.21--94.83. Their selected validation NLL is 2.641--2.712 versus Original teacher-forced dev NLL 3.235, showing that likelihood and decoded recognition competence are not interchangeable. On PHX-public, reconstruction GT scores (7.19--9.67 BLEU points) likewise remain below Original 12.78. Consequently, evaluator competence and unavailable original-training details remain confounded with evaluator identity; the experiment identifies evaluator-pipeline sensitivity, not a causal checkpoint-identity effect.

\paragraph{Competence gate and trajectory analysis.}
A checkpoint is competence-matched to the original evaluator only if both $|\Delta\text{dev BLEU-4}|\leq 1.0$ BLEU point and $|\Delta\text{dev WER}|\leq 3.0$ points (official implementations, 0--100 scale; the released evaluator's dev values are 13.38 and 83.37, the latter matching the official repository's reported 83.3702). These are prespecified operational thresholds set by the authors before any rescue run was inspected; we do not claim they are domain-standard, and all gate conclusions are reported as counts against this declared rule rather than as a principled equivalence bound. None of the fifty reconstruction-family checkpoints passes either half of this gate: dev BLEU-4 falls short by 3.4--8.9 points, and dev WER exceeds the 86.37 bound by 0.08--13.3 points. Along the dense-checkpoint training trajectories (170 saved epochs across eight reconstruction seeds (six primary with dense checkpointing, plus two extension seeds with partial trajectory)), no saved epoch reaches the $\pm 1.0$ BLEU competence band (maximum 9.73) and the best official WER across all saved epochs (86.41) still exceeds the 86.37 bound. A trajectory-based competence match is therefore infeasible without altering the training recipe.

Because checkpoint selection by validation NLL could in principle hide a competence-matched epoch, we computed the multi-objective Pareto frontier over every saved epoch of the six dense trajectories (dev NLL, decoded dev BLEU-4, decoded dev WER): no Pareto-optimal epoch under \emph{any} of the three selection objectives reaches the gate (best BLEU-selected epoch: 9.72, seed 202 epoch 50; the WER gate requires $\leq 86.37$ under the official protocol). BLEU-based selection systematically prefers later epochs (50--100) than NLL-based selection (20--35), so the two objectives diverge, but neither closes the competence gap. We also expanded the rescue search beyond learning rate: batch size $\{128, 512\}$, dropout $\{0.2\}$, weight decay $\{0\}$, label smoothing $\{0.1\}$, and a longer-budget variant (600 epochs, lr $5\times10^{-4}$), two seeds each under the same frozen architecture, manifests, and NLL selection rule (Table~\ref{tab:rescue_dev}). All twelve new runs fail the gate (dev BLEU-4 6.8--9.9); the best rescue checkpoint overall (weight-decay 0, seed 202, dev BLEU-4 9.92, official WER 86.63) remains 3.5 BLEU points below the gate floor and 0.26 WER points above the WER bound, and evaluated on PHX-public it yields GT 10.73 / PURE 10.98 (gap $+0.25$), extending the competence--gap scatter toward the gate without any reversal reappearing. Together with the eight learning-rate rescue variants reported previously (dev BLEU-4 7.9--9.7), twenty rescue runs across six hyperparameter families fail to reach competence parity. The gate conclusion is insensitive to threshold choice: across a $\pm[0.5,2.0]$ BLEU $\times$ $\pm[1.5,4.5]$ WER grid of twelve combinations, 0 of 50 trained checkpoints passes any of them. Across the fourteen checkpoints, the gap--competence Spearman correlation is $\rho=-0.34$ ($p=0.24$); with a competence range excluding the original, this test lacks power and we do not interpret it as evidence that competence is irrelevant (the dose--response analysis of \S\ref{sec:positive_evidence} addresses this directly). We frame the result as released-evaluator pipeline non-reproduction rather than checkpoint identity.

\paragraph{What the released training log establishes.}
\label{sec:original_trajectory}
The bundle's \texttt{validations.txt} gives the released evaluator's actual training trajectory (Fig.~\ref{fig:original_trajectory}): validation BLEU-4 reaches 10.0 at step 770 ($\approx$epoch 28), 12.2 at step 1{,}134 ($\approx$epoch 41), and the best 13.38 at step 1{,}820 ($\approx$epoch 66), after which it oscillates and degrades slightly (12.9 at the final step 2{,}828); dev perplexity bottoms around step 750 and then rises, so the released selection is decoded-BLEU-based, not likelihood-based. Our reconstructed trajectories are much flatter early (dev BLEU-4 8--9.7 by epoch 50--100) and plateau at $\leq$10.0 even at epoch 800. The comparison is therefore not simply an epoch-budget gap: the released run is of comparable length ($\sim$100 epochs) but keeps improving past epoch 30--60, where every constructible reconstruction plateaus. What remains unexplained by the released artifacts is the training dynamic that produced that slope --- candidate differences include the exact data pipeline/version, initialization, schedule implementation, and the trainer's dev-evaluation loop --- and we do not claim to identify it. The log also explains the released checkpoint's WER conventions: its training-time WER at the best step is 80.27 (with DEL 76.3 / INS 0.2 / SUB 3.8 decomposition), while the official jiwer implementation gives 83.37 on the same dev set and our earlier raw-split WER gave 77.49; we report the official 83.37 throughout.

\begin{figure}[t]
\centering
\includegraphics[width=0.85\textwidth]{generated_figures/original_trajectory.pdf}
\caption{Released evaluator's training trajectory from the bundle's \texttt{validations.txt} (202 validation points over 2{,}828 optimizer steps, decoded-BLEU selection): dev BLEU-4 (left axis) and dev perplexity (right axis) versus optimizer step. The released \texttt{best.ckpt} corresponds to step 1{,}820 (best dev BLEU-4 13.38); perplexity bottoms around step 750 and rises thereafter, confirming decoded-BLEU rather than likelihood selection.}
\label{fig:original_trajectory}
\end{figure}

\paragraph{Competence matching by degradation.}
\label{sec:degradation}
We cannot train a reconstruction up to the released competence (13.38), but we can degrade the released evaluator's inputs until its GT score matches the reconstruction range, and ask whether the replay gap survives (Table~\ref{tab:degradation}). Two degradation axes behave very differently: per-joint spatial jitter up to $\sigma{=}8\%$ of per-joint scale barely moves either score (GT 12.4--12.6, PURE 23.6--23.9), while temporal downsampling with interpolation back to full length steadily compresses both. At $k{=}8$ (one frame in eight kept), GT falls to 8.75 --- inside the reconstruction range (7.2--9.7) --- yet PURE still scores 11.7, above every reconstruction's \emph{clean} PURE (maximum 11.0), and the gap remains $+2.9$. The replay advantage therefore attenuates with effective input quality but is not eliminated: it is not an artifact of the released evaluator merely being a strong recognizer on clean inputs. We emphasize that temporal degradation is also a distributional shift, so this bounds, rather than identifies, the competence component of the anomaly.

\begin{table}[h]
\centering
\scriptsize\setlength{\tabcolsep}{3pt}
\caption{Degradation matching under the released evaluator (corpus sacreBLEU, BLEU points). $k$: keep every $k$-th frame, interpolate back to full length. Jitter: per-joint Gaussian noise, $\sigma$ as fraction of per-joint scale (values pooled over $\sigma\in[0,0.08]$ since the jitter axis is nearly flat). Reconstruction clean scores for comparison: GT 7.2--9.7, PURE 6.8--11.0, gap $[-1.1,+0.3]$.}
\label{tab:degradation}
\begin{tabular}{lrrrr}
\toprule
Degradation & GT & PURE & Gap & Note \\
\midrule
clean ($k{=}1$) & 12.78 & 23.79 & $+11.01$ & headline \\
$k{=}2$ & 12.77 & 23.73 & $+10.96$ & \\
$k{=}3$ & 12.20 & 22.61 & $+10.42$ & \\
$k{=}4$ & 11.77 & 20.29 & $+8.53$ & \\
$k{=}6$ & 10.39 & 15.70 & $+5.31$ & \\
$k{=}8$ & 8.75 & 11.69 & $+2.94$ & GT at reconstruction level \\
jitter $\sigma\leq0.08$ ($k{=}1$) & 12.42--12.67 & 23.78--23.89 & $+11.1$ to $+11.5$ & spatial axis nearly flat \\
\bottomrule
\end{tabular}
\end{table}

\begin{table}[h]
\centering
\small
\caption{Competence rescue on the released 515-sequence dev split: 20 runs across six hyperparameter families (two to four seeds each; per-seed values in the artifact). All fail the competence gate ($|\Delta\text{dev BLEU-4}|\leq 1.0$ and $|\Delta\text{dev WER}|\leq 3.0$ under the official protocol; the released evaluator is 13.38 / 83.37; official WER ranges: lr-rescue 87.5--93.1, expanded rescue 86.6--93.1); a long-schedule run reaches a family-maximum 10.02 and still fails. The competence gate conclusion is insensitive to threshold choice: across a $\pm[0.5,2.0]$ BLEU $\times$ $\pm[1.5,4.5]$ WER grid of twelve combinations, 0 of 50 trained checkpoints passes any of them. The multi-objective Pareto frontier over the dense trajectories independently confirms no saved epoch under any selection objective reaches the gate.}
\label{tab:rescue_dev}
\begin{tabular}{lrr}
\toprule
Variant family & dev BLEU-4 range & min $|\Delta\text{BLEU-4}|$ vs Original \\
\midrule
lr $5\times10^{-4}$ / $2\times10^{-3}$ (8 runs) & 7.9--9.7 & 3.7 \\
batch size 128 / 512 & 6.8--9.2 & 4.2 \\
dropout 0.2 & 6.9--7.5 & 5.9 \\
weight decay 0 & \textbf{9.8--9.9} & \textbf{3.5} \\
label smoothing 0.1 & 6.8--7.4 & 6.0 \\
600 epochs, lr $5\times10^{-4}$ & 8.2--8.3 & 5.1 \\
\bottomrule
\end{tabular}
\end{table}







The six primary seeds carry precomputed paired-bootstrap intervals; the eight extension seeds in Table~\ref{tab:multiseed_factorial} independently reproduce the same near-zero or negative point-estimate pattern without being folded post hoc into those six primary comparisons (seed-level 95\% prediction interval $[-1.4,+0.3]$ BLEU points).

\paragraph{Cross-metric consistency across checkpoints.}
Table~\ref{tab:cross_metric_gap} reports the PURE--GT gap on the same fixed output under each of the seven primary evaluators across decoded-text and likelihood metrics, with chrF, WER and ROUGE-L under the official implementations (\S\ref{sec:metric_protocol_alignment}). Under the original evaluator, the gap is positive and large for all decoded-text metrics; teacher-forced NLL is more negative for PURE than GT ($-0.419$, sign-flipped), so the only metric family that does not reverse is BT likelihood, consistent with the likelihood-versus-decoded-metric divergence reported by Jiang et al.~\cite{jiang2025meaningful}. Under the six reconstructions the large reversal disappears: all chrF and ROUGE-L gaps are negative ($-0.38$ to $-1.94$ and $-0.15$ to $-1.67$ respectively), all WER gaps favor GT, and sacreBLEU-4 gaps range from $-0.98$ to $+0.13$ BLEU points with only two small positive cells (seeds 101 and 606). The reversal is therefore not specific to sacreBLEU-4 or to lexical n-gram matching; it is a property of the decoded-text metrics under the original evaluator that is not reproduced under reconstructions.

\paragraph{Rank stability across evaluators on the nine-system set.}
\label{sec:rank_stability}
Beyond the four canonical systems, we score a fixed nine-system set --- GT, PT (a generation system that does not copy training poses), TN-PURE, TN-PTCOMP, LaBSE and multi-SBERT donor copy, oracle-gloss composition, UNSEEN-PURE (test-pool donors, likewise not training poses), and random donor --- under all seven primary evaluators and compute Kendall $\tau_b$, Spearman $\rho$, and pairwise flip rates. The original evaluator ranks TN-PURE first; five of six reconstructions rank GT first. Original-vs-reconstruction agreement is only moderate (mean $\tau_b=0.44$, $\rho=0.57$, flip rate 27.8\%), and reconstruction-vs-reconstruction agreement is higher but still not strong (mean $\tau_b=0.58$, $\rho=0.73$, flip rate 21.1\%). Two readings follow: rank instability is partly generic across competence levels, and the released evaluator deviates \emph{more} than the typical reconstruction pair. The set deliberately mixes constructed probes with real pipeline outputs: the official starter Progressive Transformer predictions (public), a retrieval+LM pipeline (noisy-gloss, our same-paradigm proxy), and GT recordings. The SLRTP2025 winner's checkpoints and test predictions are not public (the official repository links the top-3 teams' codebases but not their submission files), so we cannot include actual competitive entries; the noisy-gloss proxy is explicitly a stand-in, not a substitute, for them. This quantifies ranking sensitivity on our probe set; it does not refute rank stability for ordinary competitive systems, and we do not claim otherwise.

\begin{table}[t]
\centering
\scriptsize\setlength{\tabcolsep}{3pt}
\caption{PURE--GT gap across metrics and evaluators on PHX-public (641 sequences; canonical decoding) for the original evaluator and the \emph{six primary} reconstructions (of fourteen; the eight extension seeds show the same near-zero/negative pattern in Table~\ref{tab:multiseed_factorial}). chrF, WER and ROUGE-L use the official SLRTP2025 implementations (\S\ref{sec:metric_protocol_alignment}). WER and NLL are sign-flipped so positive always means PURE is better than GT. BERTScore-F1 uses \texttt{bert-base-multilingual-cased} (bert-score 0.3.12, \texttt{lang=de}, \texttt{idf=False}, no baseline rescaling); a BLEURT substitute column is in the artifact. The large reversal is consistent in direction across decoded-text metrics under the original evaluator and disappears under all six primary reconstructions; residual reconstruction signs are mixed and near zero for n-gram metrics.}
\label{tab:cross_metric_gap}
\begin{tabular}{lrrrrrrr}
\toprule
Evaluator & $\Delta$BLEU-4 & $\Delta$BLEU-1 & $\Delta$chrF & $\Delta$WER$^\dagger$ & $\Delta$NLL$^\dagger$ & $\Delta$BERTSc. \\
\midrule
Original & $+11.01$ & $+21.03$ & $+13.23$ & $+13.90$ & $-0.419$ & $+0.066$ \\
Seed 101 & $+0.08$ & $-0.15$ & $-1.31$ & $-2.95$ & $-0.219$ & $-0.005$ \\
Seed 202 & $-0.98$ & $-0.21$ & $-1.94$ & $-3.52$ & $-0.233$ & $-0.006$ \\
Seed 303 & $-0.71$ & $-0.10$ & $-1.20$ & $-1.53$ & $-0.240$ & $+0.001$ \\
Seed 404 & $-0.36$ & $-0.14$ & $-1.20$ & $-1.87$ & $-0.194$ & $-0.001$ \\
Seed 505 & $-0.89$ & $-0.17$ & $-1.78$ & $-3.41$ & $-0.268$ & $-0.004$ \\
Seed 606 & $+0.13$ & $-0.04$ & $-0.38$ & $-3.46$ & $-0.169$ & $-0.001$ \\
\bottomrule
\end{tabular}
\par\smallskip
\footnotesize $^\dagger$WER and NLL are sign-flipped: positive means PURE is better than GT (lower WER, lower NLL).
\end{table}

\paragraph{Reference-validity sensitivity: the reversal attenuates against human sign-to-text references.}
\label{sec:human_refs}
Two 2026 corpus audits bear directly on this reversal. Czehmann et al.~\cite{czehmann2026sacred} audited PHOENIX-2014T and released manual sign-to-text back-translations of the full test set (produced by one deaf fluent L2 DGS signer, with per-item confidence flags; the high-confidence subset has no second-annotator adjudication), documenting information loss, lexical errors, and translationese in the original speech-transcribed references. Alkain et al.~\cite{alkain2026critical} quantified train--test text overlap across sign-language translation corpora, showing that a small fraction of highly training-like test items disproportionately inflates corpus BLEU. We reproduce the overlap pattern locally: the 32 test items whose reference exactly matches a training text (top 5\% by max-train Jaccard) score GT 40.26 / PURE 78.73 (gap $+38.5$), while the remaining 609 items score GT 11.93 / PURE 22.20 (gap $+10.3$); the reversal persists in attenuated form even in the bottom half by overlap (GT 3.50 / PURE 11.00, gap $+7.5$). Replacing the reference texts with the human back-translations of~\cite{czehmann2026sacred} attenuates the reversal sharply (Table~\ref{tab:ref_sensitivity}): against the full human-reference set (641 items; the resource has 642 rows, one per official test ID --- our 641 reflect the same single missing released pose record as elsewhere, and the confidence=1 subset likewise drops from 462 to 461 rows) the gap is $+1.12$ (95\% CI $[+0.38,+1.91]$) versus $+11.01$ $[+9.61,+12.43]$ against the original references; the paired attenuation is $9.89$ BLEU points $[8.43,11.32]$, i.e., a remaining ratio of 10.2\% $[3.5\%,17.1\%]$. This is a \emph{descriptive contrast of two corpus gaps}, not a causal attribution: corpus BLEU is non-linear and the two gaps are measured against different reference sets. Under all six reconstruction checkpoints the gap against human references is negative ($-0.15$ to $-0.65$). Following~\cite{czehmann2026sacred}'s single- versus dual-reference protocol, scoring against \emph{both} references jointly (official + human as two references) leaves the gap at $+11.55$ (full set) and $+11.10$ (high-confidence): the replay's decode matches the original reference register far better than the recorded poses' decode matches \emph{either} reference, so the reversal is not rescued by offering the human reference as an alternative --- it is specific to the original reference frame. The original references themselves score only 11.87 BLEU against the human references, confirming the reference-reference divergence documented by~\cite{czehmann2026sacred}. Together with \S\ref{sec:factorial_corrected}, the picture is compositional rather than attributional: the large reversal requires (i) the original reference frame, and (ii) the released evaluator's familiarity with training-pool donors, which is strongest exactly where donors overlap the query (\S\ref{sec:slot_f1}).

\begin{table}[h]
\centering
\scriptsize
\caption{PURE--GT gap (corpus sacreBLEU, BLEU points) under three reference conditions for the released evaluator and the six primary reconstructions. Human references: Czehmann et al.~\cite{czehmann2026sacred} back-translations (CC BY-NC-SA 4.0). Dual: official + human as two references (sacrebleu multi-reference; closest-length reference convention). The attenuation from original to human references is descriptive, not a causal attribution.}
\label{tab:ref_sensitivity}
\begin{tabular}{lrrrr}
\toprule
Reference condition & $N$ & GT & PURE & Gap \\
\midrule
Original (speech-transcribed) & 641 & 12.78 & 23.79 & $+11.01$ $[+9.61,+12.43]$ \\
Human back-translation (full) & 641 & 5.66 & 6.78 & $+1.12$ $[+0.38,+1.91]$ \\
Human back-translation (confidence=1) & 461 & 7.35 & 8.66 & $+1.32$ \\
Dual (original + human) & 641 & 14.97 & 26.51 & $+11.55$ \\
Dual (confidence=1 subset) & 461 & 17.57 & 28.67 & $+11.10$ \\
\midrule
Human refs, reconstructions (6 seeds) & 641 & 3.49--4.62 & 3.35--4.35 & $-0.15$ to $-0.65$ \\
\bottomrule
\end{tabular}
\end{table}

\paragraph{Transcript-slot overlap diagnostic (not an independent semantic criterion).}
\label{sec:slot_f1}
Decoded-text metrics all consume the same BT decode, and the slot diagnostic in this paragraph shares that input: it is \emph{not} an independent judgment of sign-language semantics --- it observes no spatial grammar, indexation, or non-manual markers --- it is a reference-anchored text diagnostic with a different computational form than n-gram metrics. We extract factual slots from the reference and hypothesis texts with a frozen rule-based German weather extractor (number/temperature, time, place, and weather-event families; lexicon, normalization rules, and per-family precision/recall in the artifact, in the spirit of slot-level semantic checks such as SLU-2K~\cite{slu2k}) and compute multiset slot-F1. Under the original evaluator, the replay's decoded text attains higher slot-F1 than the reference poses' own decode (0.526 vs 0.354 micro-F1 over all slot families), so the reversal extends beyond n-gram form to factual content \emph{as read by this evaluator}. The donor \emph{transcript} itself, without any evaluator, attains slot-F1 0.575 against the query reference: the replay decode reaches 91\% of this donor-transcript \emph{baseline} (not a ceiling --- an evaluator could in principle paraphrase beyond the transcript), consistent with the released evaluator passing through donor content it was trained on. Under all six reconstructions the ordering flips (GT 0.206--0.239 vs PURE 0.177--0.234). To validate the extractor we audited a frozen uniform sample of 50 references (seed 20260729): per-family precision 0.91--1.00, recall 0.75--1.00 (overall precision 0.97, recall 0.86); errors are dominated by recall misses from inflections, compounds, and lexicon gaps. A further stratified audit of 50 decoded hypotheses (25 GT-decode, 25 PURE-decode, seed 20260730) shows extractor behavior is \emph{not} system-independent (Table~\ref{tab:slot_audit}): precision is nearly equal (0.96 vs 0.97) but recall is lower on PURE-decode (0.83) than on GT-decode (0.91), because donor-transcript-style outputs use more inflected verbs and compounds (\emph{regnet}, \emph{unwetterwarnungen}) that the lexicon misses. We therefore do not claim cross-system unbiasedness; since precision is nearly equal and recall is lower on PURE, the measured PURE slot-F1 is attenuated more than GT's, so the PURE$>$GT contrast is if anything underestimated. The audit used a single auditor; inter-annotator agreement is unavailable and is listed as a limitation.

\begin{table}[h]
\centering
\scriptsize\setlength{\tabcolsep}{3pt}
\caption{Slot-extractor audit by output type (50 items each; single auditor). Sentence-level micro aggregates. Reference-side audit: overall P 0.97 / R 0.86.}
\label{tab:slot_audit}
\begin{tabular}{lrrrr}
\toprule
Output type & Precision & Recall & F1 & Dominant error types \\
\midrule
References ($n=50$) & 0.97 & 0.86 & 0.91 & inflections, compounds, lexicon gaps \\
GT-decode ($n=25$) & 0.96 & 0.91 & 0.93 & pressure-system names (tief/hoch), compounds \\
PURE-decode ($n=25$) & 0.97 & 0.83 & 0.89 & inflected verbs (regnet), unwetterwarnungen, compounds \\
\bottomrule
\end{tabular}
\end{table}

The seven-evaluator ranking matrix shows the same association alongside score compression: its rows are deliberately constructed probes (replay, composition, weak generation, random controls) rather than a representative sample of competitive SLP systems. Decoder settings are fixed (beam 3, length penalty $-1$, max length 30; sacreBLEU signature \texttt{BLEU|nrefs:1|case:mixed|eff:no|tok:13a|smooth:exp|version:2.5.1}); the canonical grid has 60 cells (15 evaluators $\times$ 4 canonical systems, distinct from the 63-cell nine-system rank matrix of \S\ref{sec:rank_stability}), with empty-hypothesis rate zero throughout, and the reversal is expressed through much higher n-gram precision for PURE (P1 60.6 vs 37.0; full decomposition in the artifact). Text-nearest retrieval ranks first and the reference system fourth under the original evaluator; the reference system ranks first in four of six reconstructions, which also score the reference poses lower (8--10 vs 12.78 BLEU points) --- evaluator identity, training randomness, optimization trajectory, and competence remain confounded.



\paragraph{Template-density and holdout boundary.}
Within PHX-public, retrieval exceeded the recorded test poses even in the low-density stratum ($+10.7$ BLEU points), so template density moderates magnitude rather than gating the reversal; template-disjoint and BT-retrained holdouts and donor-familiarity cross-fitting do not reproduce it (Appendix~\ref{sec:appendix_template_boundary}).

\subsection{Donor-origin contrast retrieval characterizes the reversal's co-occurrence with training-pool donors}
\label{sec:exposure_results}

If the released-evaluator reversal requires exposure to training-pool donors, then restricting retrieval to poses the released evaluator has not seen should remove or invert the advantage. We test this directly using the exposure-symmetric systems of \S\ref{sec:exposure_systems}, holding the evaluator and query set fixed.

Under the original SLRTP2025 evaluator, UNSEEN-PURE-v1 scores 8.86 sacreBLEU (0--100 scale) versus 12.78 for the recorded test poses, a gap of $-3.92$ BLEU points: the retrieval advantage observed for SEEN-PURE-v1 ($+11.01$ BLEU points) does not survive retrieval-pool substitution. The donor-pool substitution estimand
\[
(\text{SEEN-PURE}-\text{GT})-(\text{UNSEEN-PURE}-\text{GT})=+14.93\text{ BLEU points}
\]
is large and signed, but is a \emph{donor-pool substitution effect}, not a pure exposure effect: SEEN-PURE-v1 and UNSEEN-PURE-v1 differ simultaneously in pool size (7060 versus 640), median source-text Jaccard (0.50 versus 0.37), donor reuse concentration (Gini 0.11 versus 0.29), and evaluator training-pool exposure. Cross-fit controls reproduce the same direction (CROSSFIT-A 8.74, CROSSFIT-B 9.13; pooled CROSSFIT-PURE minus fold-local-GT is $-3.8$ BLEU points). Table~\ref{tab:exposure_estimands} reports the per-evaluator cells. (The 302-sequence UNSEEN-MATCHED-v1 subset scores 8.86 under the original evaluator, a display-precision coincidence with the full-support 8.86 that we note only for completeness.)

\begin{table}[t]
\centering
\scriptsize\setlength{\tabcolsep}{3pt}
\caption{Donor-origin contrast estimands (sacreBLEU, BLEU points; gaps vs local GT). SEEN = 7,060 training poses; UNSEEN = 640 other test poses; CROSSFIT = hash-split folds. Per-seed cells in the artifact.}
\label{tab:exposure_estimands}
\begin{tabular}{lrrrrr}
\toprule
Evaluator & SEEN-PURE & UNSEEN-PURE & SEEN$-$GT & UNSEEN$-$GT & Estimand \\
\midrule
Original & 23.79 & 8.86 & $+11.01$ & $-3.92$ & $\mathbf{+14.93}$ \\
Seeds (range) & 6.83--9.09 & 6.59--7.98 & $-1.10$ to $+0.13$ & $-1.94$ to $-0.27$ & $-0.15$ to $+1.22$ \\
CROSSFIT pooled & --- & 8.74 / 9.13 & --- & $-3.8$ & --- \\
\bottomrule
\end{tabular}
\end{table}

Three cautions bound this result: the pools differ in similarity (median Jaccard 0.50 vs 0.37, though donor-exclusion controls show the SEEN-PURE reversal surviving high-overlap exclusion to $\tau{=}0.5$, \S\ref{sec:appendix_public_donor_exclusion}), in reuse concentration (561 vs 358 unique donors, Gini 0.11 vs 0.29), and the cross-fit pooled estimand uses fold-local GT re-aggregation.

\paragraph{Donor-pool substitution: path-dependent descriptive contrasts.}
\label{sec:decomposition}
We decompose the total substitution estimand only descriptively, along explicitly labelled telescoping paths. With SEEN-RAND640-MATCHED-v1 added (\S\ref{sec:exposure_systems}), the systems form a $2\times2$ within-train design (pool size $\{$7060, 640$\}$ $\times$ selection $\{$max-Jaccard, Jaccard-matched$\}$) plus the test-pool cell: S1 = SEEN-PURE (7060, max), S2 = SEEN-RAND640 (640, max), S3 = SEEN-MATCHED (7060, match), S4 = SEEN-RAND640-MATCHED (640, match), U = UNSEEN-PURE (test 640, max).
Because corpus BLEU is non-linear in the support, all contrasts use the \emph{common} 622-query support (cells S1~$=23.6$, S2~$=16.5$, S3~$=17.3$, S4~$=14.7$, U~$=8.9$ under the original evaluator), where (A)~$=$ S1$-$U telescopes exactly along any ordering. The prespecified canonical path (S1$\to$S2$\to$S3$\to$U) gives size $+7.2$ $[+5.9,+8.4]$, similarity $-0.8$ $[-1.8,+0.2]$, origin $+8.4$ $[+7.3,+9.4]$; the two alternative admissible orderings give materially different attributions (Table~\ref{tab:path_sensitivity}): the similarity/selection term ranges from $-0.8$ to $+6.4$ BLEU points and the origin term from $+8.4$ to $+5.8$ across orderings, because size and selection interact strongly. Shapley-style averages over the admissible orderings are: size $+4.9$ $[+4.0,+5.8]$, selection $+4.1$ $[+3.3,+4.9]$, origin $+5.8$ $[+4.7,+6.8]$. We report all of these as \emph{path-dependent descriptive contrasts}, not mechanism decompositions; the better-balanced factorial (\S\ref{sec:factorial_corrected}) estimates the pool-origin main effect at $+2.1$, below every path-based origin estimate. Under the six reconstructions all pool-origin contrasts lie in $[-0.6,+0.6]$ with CIs including zero. Pool-to-pool variability in the size contrast is small (20 independent shared 640-donor pools: mean $16.6$, SD $0.5$, containing the per-query S2 score of $16.5$).

\begin{table}[h]
\centering
\scriptsize
\caption{Path sensitivity of the donor-pool substitution decomposition under the original evaluator (common 622-query support; 10{,}000-resample query-index bootstrap CIs). S1=SEEN-PURE (7{,}060-pool, max-Jaccard), S2=SEEN-RAND640 (640-pool, max), S3=SEEN-MATCHED (7{,}060-pool, Jaccard-matched), S4=SEEN-RAND640-MATCHED (640-pool, matched), U=UNSEEN-PURE (test pool). Every ordering telescopes to the same total (A)~$=$ S1$-$U~$=+14.7$ $[+13.3,+16.2]$ exactly, but the attribution to individual factors varies by up to 7 BLEU points across orderings; all contrasts are descriptive.}
\label{tab:path_sensitivity}
\begin{tabular}{llrrr}
\toprule
Ordering & Step & Size contrast & Selection contrast & Origin contrast \\
\midrule
Canonical (prespecified) & S1$\to$S2$\to$S3$\to$U & $+7.2$ $[+5.9,+8.4]$ & $-0.8$ $[-1.8,+0.2]$ & $+8.4$ $[+7.3,+9.4]$ \\
Alternative 1 & S1$\to$S2$\to$S4$\to$U & $+7.2$ $[+5.9,+8.4]$ & $+1.8$ $[+1.1,+2.5]$ & $+5.8$ $[+4.7,+6.8]$ \\
Alternative 2 & S1$\to$S3$\to$S4$\to$U & $+2.6$ $[+1.7,+3.5]$ & $+6.4$ $[+5.1,+7.7]$ & $+5.8$ $[+4.7,+6.8]$ \\
\midrule
Shapley average & over orderings & $+4.9$ $[+4.0,+5.8]$ & $+4.1$ $[+3.3,+4.9]$ & $+5.8$ $[+4.7,+6.8]$ \\
\bottomrule
\end{tabular}
\end{table}

\paragraph{Uncertainty: bootstrap family.}
All decomposition CIs use 10{,}000 resamples on the full 622-query support. We report four resampling designs: a query-index paired bootstrap (which jointly resamples each query's deterministically linked donor IDs); three separate one-way cluster bootstraps for the origin contrast (template-family, SEEN-donor, and UNSEEN-donor clusters); and a joint pigeonhole multiway replicate that resamples all three dimensions per replicate. The origin-contrast CIs are $[+7.2,+9.5]$ / $[+7.1,+9.7]$ / $[+7.1,+9.7]$ / $[+7.7,+9.1]$ respectively, all matching the query-index CI $[+7.3,+9.4]$ to within $\pm0.3$, so dependence through shared donors is negligible. Fuzzy template families (slot-masked normalized references) give 589 clusters over 622 queries (mean size 1.06), so template sharing is rare and template-cluster CIs match sequence CIs. The headline gap uses the paired permutation test: 0 of 10{,}000 exceedances, $p\leq10^{-4}$ (95\% binomial CI on the rate $[0,0.00037]$).

\paragraph{The 605-query factorial: full $2\times2$ report.}
\label{sec:factorial_corrected}
The complementary four-cell factorial matches evaluator-seen training donors and evaluator-unseen test donors at high and low LaBSE similarity under predeclared tolerances (Table~\ref{tab:factorial_design}), retaining 605 of 641 targets (36 excluded for \texttt{no\_joint\_match\_within\_predeclared\_tolerances}; selection frozen before scoring). All six conditions are scored at the canonical 12.5\,fps --- the external \texttt{[::2]} transform is required because the released BT model applies no internal subsampling (\S\ref{sec:method}). Table~\ref{tab:factorial_full} reports the complete $2\times2$ cells, both main effects, and the interaction with 10{,}000-resample CIs and Holm correction. Under the original evaluator the seen (pool-origin) main effect is $+2.08$ $[+1.56,+2.57]$ (Holm $p=0.0003$), the similarity main effect $+6.96$ $[+5.93,+8.08]$, and the seen$\times$similarity interaction $+3.70$ $[+2.63,+4.79]$ (Holm $p=0.0003$); NLL shows the same pattern ($-0.128$ and $-0.124$, Holm $p=0.0045$). Across reconstructions the seen main effect is smaller ($+0.51$ $[+0.25,+0.77]$) and the interaction is compatible with zero ($-0.39$ $[-0.92,+0.14]$). Two conclusions follow, with one caveat: the balanced pool-origin estimate ($+2.1$) is much smaller than any path-based origin contrast ($+5.8$--$+8.4$), so the Jaccard-only path contrasts inflate the pool-origin component, and we take the factorial value as the more credible descriptive one; and the released evaluator's seen advantage \emph{grows with donor--query similarity} while reconstruction interactions are null, concentrating the anomaly in the high-similarity regime. The caveat: even within strata, seen donors are systematically more similar to their queries than unseen donors (similarity SMD $+0.41$/$+0.38$; the 7{,}060-pose train pool offers closer matches), so the seen main effect partly inherits that residual imbalance and remains a train-pool vs test-pool provenance contrast, not an identified training-exposure effect.

\begin{table}[h]
\centering
\small
\caption{Factorial design and balance report (605 retained targets; 36 excluded for no joint match within tolerances). High: seen/unseen donor pair maximizing LaBSE cosine with $|$sim(seen)-sim(unseen)$|\leq0.08$; low: similarity $\leq$ high$-0.12$, same cross-pool tolerance. Calipers: candidate--target relative duration difference $\leq0.35$, four-cell duration range $\leq0.30$ (low-to-high $\leq0.20$); word-count difference $\leq8$, four-cell range $\leq6$ (low-to-high $\leq4$); all four cells share the query's signer. Selection: high pair first (maximize summed similarity, then cross-pool $|\Delta$sim$|$, then lexical IDs), then low pair; deterministic tie-breaks. SMD = standardized mean difference, seen vs unseen within stratum. Cross-cell donor collisions were forbidden by construction; 1{,}519 unique donors fill the 2{,}420 cell slots (max reuse 10).}
\label{tab:factorial_design}
\begin{tabular}{lrrrr}
\toprule
Quantity & seen\_high & seen\_low & unseen\_high & unseen\_low \\
\midrule
LaBSE similarity (mean) & 0.813 & 0.654 & 0.776 & 0.622 \\
Duration (frames, mean) & 100.2 & 99.5 & 99.2 & 98.8 \\
Word count (mean) & 13.3 & 13.6 & 13.3 & 13.4 \\
Same-signer rate & 100\% & 100\% & 100\% & 100\% \\
\midrule
Similarity SMD (seen $-$ unseen) & \multicolumn{2}{r}{$+0.41$ (high stratum)} & \multicolumn{2}{r}{$+0.38$ (low stratum)} \\
Duration / word-count SMD & \multicolumn{4}{l}{$|SMD|\leq0.04$ in both strata} \\
\bottomrule
\end{tabular}
\end{table}

\begin{table}[h]
\centering
\scriptsize\setlength{\tabcolsep}{3pt}
\caption{Full $2\times2$ exposure factorial at the canonical 12.5\,fps protocol (605 targets; corpus sacreBLEU in BLEU points; NLL per token). Cells: seen/unseen = evaluator-training-pool vs public-test-pool donors; high/low = LaBSE similarity strata, matched within predeclared duration, word-count, and same-signer tolerances. CIs: 10{,}000-resample bootstrap (original) and seed-level $t$ across six reconstructions; Holm correction across the three contrasts.}
\label{tab:factorial_full}
\begin{tabular}{llrrrrrrr}
\toprule
Metric & Evaluator & seen\_hi & seen\_lo & unseen\_hi & unseen\_lo & seen main & sim.\ main & inter. \\
\midrule
BLEU & Original & 11.12 & 2.31 & 7.19 & 2.08 & $\mathbf{+2.08}$ & $+6.96$ & $\mathbf{+3.70}$ \\
 &  &  &  &  &  & $[+1.6,+2.6]$ & $[+5.9,+8.1]$ & $[+2.6,+4.8]$ \\
BLEU & Seeds (mean) &  &  &  &  & $+0.51$ & $+4.06$ & $-0.39$ \\
 &  &  &  &  &  & $[+0.3,+0.8]$ & $[+3.6,+4.5]$ & $[-0.9,+0.1]$ \\
\midrule
NLL & Original & 3.74 & 4.30 & 3.93 & 4.37 & $-0.128$ & $-0.498$ & $-0.124$ \\
 &  &  &  &  &  & $[-.17,-.09]$ & $[-.55,-.45]$ & $[-.21,-.04]$ \\
NLL & Seeds (mean) &  &  &  &  & $-0.064$ & $-0.179$ & $-0.004$ \\
 &  &  &  &  &  & $[-.08,-.05]$ & $[-.20,-.16]$ & $[-.01,+.00]$ \\
\bottomrule
\end{tabular}
\end{table}



\paragraph{Matched-support balance and overlap.}
The three matched supports are: SEEN-PURE-MATCHED-v1 (622 queries, Jaccard-matched train donors within tolerance 0.1 of each UNSEEN donor), UNSEEN-PURE-v1 (641 queries, retrieval from the 640 other test poses), and UNSEEN-MATCHED-v1 (302 queries, subset admitting a SEEN donor within Jaccard and 1\,s duration tolerance). Their Jaccard distributions match in central tendency (means 0.411/0.424/0.432, medians 0.36--0.37); the canonical SEEN-PURE-v1 is higher (0.535/0.500) with 32 exact normalized-text matches (5.0\%) and 33 token-multiset-exact pairs (5.1\%), while the three matched supports have 0 exact normalized-text matches by construction and 9 token-multiset-exact pairs (identical token multisets under word-order permutation, e.g., \emph{"guten abend liebe zuschauer"} vs \emph{"liebe zuschauer guten abend"}). Excluding those 9 queries leaves the pool-origin residual unchanged ($+8.4$, $n{=}613$). Donor reuse: SEEN-PURE-v1 uses 561 unique donors (87.5\%, max reuse 7, Gini 0.11); UNSEEN-PURE-v1 uses 358 (55.9\%, max reuse 8, Gini 0.29).

\paragraph{Multivariate balance, caliper sensitivity, and signer-strict matching.}
\label{sec:balance}
Beyond Jaccard, Table~\ref{tab:multivariate_balance} reports standardized mean differences (SMD) between the SEEN-PURE-MATCHED and UNSEEN-PURE donors on the 622-query support for six donor covariates: source-text Jaccard, LaBSE cosine similarity, donor duration, donor word count, donor template density (gloss-bigram Jaccard to the nearest train neighbour), and same-signer indicator. All $|SMD|<0.12$, so the two donor groups are marginally well balanced on observed covariates. Per-query pairing is looser: only 57\% of pairs have $|\Delta$duration$|\leq1$\,s and 69\% have $|\Delta$words$|\leq2$, and both donors share the query's signer in only 4.7\% of pairs. Pose missingness is not observable in the released tensors (no NaN or all-zero joints in any released pose record), so balance can only be assessed on the observed covariates. Caliper sensitivity of the pool-origin residual (SEEN-MATCHED $-$ UNSEEN-PURE on the respective support, original evaluator) is small: tolerance 0.05 ($n=597$) gives $+8.2$ $[+7.2,+9.3]$, canonical tolerance 0.10 ($n=622$) $+8.4$, and tolerance 0.20 ($n=636$) $+8.6$ $[+7.5,+9.6]$ BLEU points. The \texttt{speaker} field is populated for all 7{,}060/515/641 released train/dev/test records, so a signer-strict sensitivity check is feasible: enforcing same-signer matching (tolerance 0.10, $n=580$ retained) attenuates the residual to $+7.0$ $[+5.9,+8.1]$ BLEU points --- the sign and significance are unchanged, so the pool-origin residual is not primarily a signer-composition artifact, though signer independence is a broader benchmark concern~\cite{artiaga2025signer}. As a stronger balance check, inverse propensity weighting (logistic propensity of train-pool origin given LaBSE, Jaccard, duration, and word count; stabilized IPW; propensity refit in each of 1{,}000 bootstrap replicates) drives all four covariate SMDs to 0.00 and leaves the residual at $+8.5$ $[+7.1,+10.1]$ --- the path-D estimate is therefore robust to exact observed-covariate balancing, though it remains an association, not an identified exposure effect (\S\ref{sec:factorial_corrected} gives the smaller factorial estimate).

\begin{table}[h]
\centering
\scriptsize\setlength{\tabcolsep}{3pt}
\caption{Multivariate donor balance on the 622-query common support: SEEN-PURE-MATCHED (train-pool) vs UNSEEN-PURE (test-pool) donors. SMD = standardized mean difference (unpooled). All $|SMD|<0.12$. Per-query overlap columns show the fraction of pairs within the stated tolerance. Pose missingness is not observable in the released tensors.}
\label{tab:multivariate_balance}
\begin{tabular}{lrrrr}
\toprule
Covariate & SEEN-MATCHED mean & UNSEEN mean & SMD & Per-query overlap \\
\midrule
Source-text Jaccard & 0.411 & 0.415 & $-0.021$ & $|{\Delta}J|{\leq}0.05$: 96\%; ${\leq}0.10$: 100\% \\
LaBSE cosine & 0.782 & 0.774 & $+0.068$ & --- \\
Duration (s) & 3.76 & 3.74 & $+0.014$ & $|{\Delta}|\leq1$\,s: 57\% \\
Word count & 12.29 & 12.31 & $-0.003$ & $|{\Delta}|\leq2$: 69\% \\
Template density & 0.310 & 0.326 & $-0.075$ & --- \\
Same signer as query & 0.183 & 0.228 & $-0.111$ & both same-signer: 4.7\% \\
\bottomrule
\end{tabular}
\end{table}



\subsection{Where the released pipeline differs: five related diagnostic views}
\label{sec:positive_evidence}
The non-reproduction result above is negative: no recipe-conditioned reconstruction shows the stress-test susceptibility. We now turn that negative statement into a positive characterization of \emph{what} distinguishes the released evaluator, through five related diagnostic views that do not rely on the competence-confounded gap comparison alone (Fig.~\ref{fig:dose_response}). These views are \emph{not} independent confirmations: they reuse the same released checkpoint, the same PHX-public data, and related donor constructions, and are best read as triangulation on one signature. The frozen list, fixed before the big-arch, long-schedule, and confirmation runs were executed, is: (i) gap--competence dose--response; (ii) donor-transcript pass-through; (iii) readout-with-generalization (in-sample vs dev/test decoding); (iv) temporal-order perturbation sensitivity; (v) constructible-protocol coverage (rescue, config-faithful, big-arch, long-schedule). The frozen-protocol confirmation at the end of this section re-runs views (i)--(iii) on two fresh seeds.

\begin{figure}[t]
\centering
\includegraphics[width=\textwidth]{generated_figures/dose_response.pdf}
\caption{Positive characterization of the released evaluator. (a) PURE--GT gap versus dev competence across 48 trained checkpoints (14 reconstructions, 12 rescue variants, 4 ladder, 4 config-faithful, 4 big-arch, 2 confirmation seeds, plus 8 lr-rescue evaluated for the gate only): all lie flat near zero; only the released evaluator (star) shows the reversal. (b) Donor-content pass-through ratio (sentence-level mean BLEU to donor transcript over sentence-level mean BLEU to query) versus dev competence: the family trends mildly upward with competence (1.05--1.31); the released evaluator sits at 3.09, a 2.4$\times$ outlier.}
\label{fig:dose_response}
\end{figure}

\paragraph{Competence dose--response is flat across the achievable range.}
We trained the documented recipe on seeded train-pool subsamples (12.5\%/25\%/50\%/75\%; 882--5{,}295 sequences) to fill the low-competence region. Together with the fourteen reconstructions and the twenty rescue runs, the competence-span set covers dev BLEU-4 $4.5$--$10.0$ across 48 checkpoints in seven families (Table~\ref{tab:checkpoint_registry}); the gate is evaluated on the 50 checkpoints of Table~\ref{tab:checkpoint_registry} that use the released train pool. The PURE--GT gap stays within $[-1.1,+0.3]$ BLEU points across this entire range (ladder: $+0.05$, $-0.20$, $+0.24$, $+0.19$; reconstructions: $-1.10$ to $+0.13$; best rescue wd0: $+0.25$), with no monotone competence trend (Fig.~\ref{fig:dose_response}a). The released evaluator, at dev BLEU-4 $13.4$, sits $+11$ BLEU points above that flat response. We stress that this excludes only the observed range: within dev BLEU-4 4.5--10.0 there is no gap trend, and we make no claim about the unobserved 10.0--13.4 region, where non-linear behavior remains possible.

\paragraph{Donor-content pass-through is a 2.4$\times$ outlier.}
For every checkpoint we decoded TN-PURE-v1 and compared the decoded text against the donor's \emph{own transcript} versus the query reference (per-item sentence BLEU, reported as sentence-level means; corpus BLEU elsewhere is always the aggregate). The released evaluator reproduces donor transcripts at mean sentence-BLEU 0.795 versus 0.257 against the query (pass-through ratio 3.09, 95\% bootstrap CI $[2.89,3.33]$; a chrF-based corroboration gives ratio 1.71 $[1.66,1.77]$). All fourteen reconstructions, all ladder checkpoints, all big-arch variants, both confirmation seeds, and both high-competence rescue runs fall in a narrow band of 1.05--1.31 with a mild competence trend (0.076--0.179 absolute donor BLEU). Corpus-level aggregation tells the same story with absolute scores: the released evaluator decodes donor poses to text at 81.2 corpus BLEU against the donor transcript versus 23.8 against the query reference (ratio 3.42), while reconstructions score 10.8--12.8 versus 8.2--9.0 (ratios 1.31--1.43). An exploratory linear fit of pass-through ratio on dev BLEU over the family ($n{=}23$ checkpoints; these are not IID but share data, code, and hyperparameters, so the interval is a visual guide, not a calibrated inference) predicts 1.30 with a nominal 95\% interval $[1.20,1.40]$ at dev BLEU 13.38; the observed 3.09 is far outside it, and the median per-item ratio (3.63 vs 1.00--1.10) shows the result is not driven by a few extreme items. (Ratios use sentence-level mean BLEU on the 0--1 scale with exp smoothing; corpus ratios use aggregate BLEU on the 0--100 scale.)

\paragraph{Readout \emph{with} generalization: the released evaluator's distinctive combination.}
\label{sec:insample}
Decoding the full 7{,}060-item training pool under the released evaluator gives free-decode BLEU 78.4, exact-match rate 69.4\%, teacher-forced top-1 token accuracy 96.5\% (train NLL 0.17), versus 13.38/12.78 free BLEU and $\approx$3\% exact match on dev/test --- it \emph{both} reads out its training pool \emph{and} generalizes to unseen data. The recipe-conditioned family separates these two regimes. On the one hand, validation-selected checkpoints (including the most overfitting-prone rescue variant, weight-decay 0) show free-decode train BLEU 11.5 $\approx$ dev 9.9: no in-sample advantage at all, and on the 525-sequence seen holdout they reach only 4.9--10.4 where the released evaluator reaches 77.0 (Table~\ref{tab:holdout_series}). On the other hand, if we let the same rescue run overfit to its final epoch (train cross-entropy 0.07/token), the training pool \emph{is} read out perfectly (free-decode train BLEU 99.9, exact match 99.5\%, NLL 0.002) --- but dev likelihood collapses (4.90), so this readout is useless as a recognizer. The family therefore \emph{can} achieve full readout, but only at the price of generalization collapse; the released checkpoint uniquely combines strong readout with strong generalization. (All donor tensors in our systems come from the released \texttt{train.pt} named as the training data in the released config; the released checkpoint's actual training inputs are unpublished, so byte-level identity of its training inputs cannot be verified --- we report the signature, not the proof, of readout.)

\paragraph{The decoder is sensitive to temporally ordered donor motion.}
A perturbation battery on the fixed GT and PURE outputs (Table~\ref{tab:perturbation}) shows the decoder is sensitive to temporally ordered donor motion: reversing all frames collapses PURE from 23.79 to 7.72 under the released evaluator (differential $16.1$ BLEU points, 95\% CI $[14.7,17.5]$; GT falls only to 7.25, differential $5.5$ $[4.6,6.5]$). We note this establishes temporal-order sensitivity, not semantic understanding: reversal, shuffling, and joint masking are also out-of-distribution corruptions, so the differentials bound how much of the score relies on ordered content rather than proving correct linguistic decoding. Window-shuffling within one-second windows costs PURE $5.4$ $[4.4,6.6]$ (GT $0.8$); 5\% spatial jitter is free ($-0.1$ $[-0.4,+0.0]$); masking hand or face joints collapses both systems (differentials $19.7$--$22.5$), so both channels carry the decoded content. The best rescue checkpoint, on identical inputs, loses only $-4.3$ under reversal for PURE and $-4.5$ for GT: it is far less sensitive to donor temporal order in the first place.

\begin{table}[h]
\centering
\scriptsize\setlength{\tabcolsep}{3pt}
\caption{Pose-perturbation battery (corpus sacreBLEU, BLEU points) under the released evaluator and the best rescue checkpoint (wd0, seed 202). Differentials with 95\% bootstrap CIs (original evaluator): reversal PURE $16.1$ $[14.7,17.5]$ vs GT $5.5$ $[4.6,6.5]$; window shuffle PURE $5.4$ $[4.4,6.6]$ vs GT $0.8$; jitter is free. Window-shuffle values use one seeded realization (a second realization gives 19.58 vs 18.35 for PURE, same conclusion).}
\label{tab:perturbation}
\begin{tabular}{llrrrrrr}
\toprule
Evaluator & System & identity & reversal & window shuffle & jitter 5\% & hand mask & face mask \\
\midrule
Original & GT-v1 & 12.78 & 7.25 & 11.97 & 12.51 & 1.35 & 3.15 \\
Original & TN-PURE-v1 & 23.79 & 7.72 & 18.35 & 23.91 & 1.30 & 4.06 \\
wd0 (seed 202) & GT-v1 & 10.73 & 6.26 & 10.54 & 10.60 & 0.63 & 1.58 \\
wd0 (seed 202) & TN-PURE-v1 & 10.98 & 6.66 & 9.65 & 10.92 & 0.58 & 1.43 \\
\bottomrule
\end{tabular}
\end{table}

\paragraph{Epoch-validation approximation of the released configuration.}
\label{sec:config_faithful_results}
The released \texttt{config.yaml} specifies \texttt{validation\_freq=14}. JoeyNMT 1.2 interprets this as 14 optimizer \emph{steps} (not epochs): at batch 256 over 7{,}060 sequences ($\approx$28 steps/epoch) this means validating approximately twice per epoch, matching the \texttt{validations.txt} log of 202 validations over 2{,}828 steps. Our initial implementation validated every 14 \emph{epochs} (an epoch-validation approximation of the configuration, not a faithful reproduction); we report those four runs here and note the protocol deviation explicitly. A step-faithful re-run (validating every 14 optimizer steps) is reported in \S\ref{sec:stepfaithful}; its interim result is that the reconstruction's dev BLEU-4 still plateaus below the original's trajectory at matched step counts. We trained four seeds under the epoch-validation approximation (Table~\ref{tab:config_faithful}). All four early-stop below the competence gate (best dev BLEU-4 8.6--9.5, vs the gate floor 12.38), all four show negative PURE--GT gaps on PHX-public ($-1.1$ to $-0.2$), and neither their donor pass-through (1.21--1.24) nor their seen-holdout readout (7.4--7.6) exceeds the recipe-conditioned family range. One seed (101) comes closest on the WER half but still fails it (86.45 vs the 86.37 bound) while remaining far from the BLEU half. The config-faithful protocol therefore does not reproduce the released checkpoint's competence, its readout signature, or the reversal. We additionally scaled capacity up (four larger variants: 4--6 layers, 384--512 hidden, 1{,}024--2{,}048 ff): these overfit faster (best epoch 15--20), reach only dev BLEU 6.8--8.0, and show family-range pass-through (1.15--1.24) and holdout scores (5.9--7.7) --- capacity alone does not recover the released evaluator's behavior either. A long-schedule continuation of the config-faithful protocol (800 epochs, no early stop, two seeds) pushes competence to a new family maximum (best dev BLEU-4 10.02 and 9.85 at epochs 252/294; GT on PHX-public 10.14) yet still shows no reversal (gap $-0.45$, holdout 8.3), and post-best trajectories degrade (dev BLEU-4 7.6--8.1 by epoch 800): competence improvement does not bring the reversal back. Within the space of publicly constructible protocols, the released pipeline's behavior is not reproducible, pointing to unpublished training details (data handling, trajectory, or selection) as the remaining difference; positively, the released checkpoint is distinguished by the readout-with-generalization signature of \S\ref{sec:insample} that no constructible protocol exhibits.

\begin{table}[h]
\centering
\scriptsize\setlength{\tabcolsep}{3pt}
\caption{Epoch-validation approximation (validate every 14 epochs, not steps; see text) (up to 3{,}000 epochs, dev-BLEU early stopping, patience 5, validation every 14 epochs, plateau $\times0.8$). All four early-stop below the competence gate (dev BLEU-4 floor 12.4); none shows the reversal, elevated pass-through, or the readout signature. Pass-through = mean sentence-BLEU(decode, donor transcript) / mean sentence-BLEU(decode, query); holdout = GT BLEU on the seen-in-training 525-sequence holdout.}
\label{tab:config_faithful}
\begin{tabular}{lrrrrrr}
\toprule
Seed & ep. & dev BLEU & dev WER & gap & pass-thr. & holdout \\
\midrule
101 & 98 & 8.91 & 86.45 & $-1.02$ & 1.24 & 7.60 \\
202 & 42 & 8.59 & 87.96 & $-1.06$ & 1.22 & --- \\
303 & 210 & 8.67 & 88.77 & $-0.23$ & 1.21 & --- \\
404 & 140 & 9.51 & 86.57 & $-0.37$ & 1.24 & 7.43 \\
\midrule
Released evaluator & --- & 13.38 & 83.37 & $+11.01$ & 3.09 & 77.0 \\
Family range (48 ckpts) & --- & 4.5--10.0 & 86.5--96.7 & $-1.10$ to $+0.25$ & 1.05--1.31 & 4.9--10.4 \\
\bottomrule
\end{tabular}
\end{table}

\paragraph{Synthesis.}
Across the five related diagnostic views within the observed competence range (dev BLEU-4 4.5--10.0), no reconstruction shows the reversal: it shows a readout-with-generalization signature (\S\ref{sec:insample}) and reproduces donor transcripts at 2.4$\times$ the family trend, and its replay advantage tracks that readout. The recipe-conditioned family, trained on identical data, shows none of these properties at any reachable competence. The reversal is therefore best described as a property of the released checkpoint and its data pipeline, co-occurring with unusually strong train-pool readout under the original reference frame --- the specific mechanism is not identifiable --- and not a property of BT evaluation generally or evidence that a randomly reseeded replica would behave alike.

\paragraph{Step-faithful re-run (correct protocol).}
\label{sec:stepfaithful}
Because the epoch-validation approximation of \S\ref{sec:config_faithful_results} deviates from the released configuration's step-based validation frequency by a factor of $\approx$28, we additionally ran two seeds under the corrected protocol (validate every 14 optimizer steps, patience 5 validations, decrease\_factor 0.8, learning\_rate\_min $10^{-8}$, stop at LR minimum). Both step-faithful runs completed (stopped at LR minimum, $10^{-8}$): seed 101 reached best dev BLEU-4 7.03 at step 1{,}470 (stopped at step 4{,}200, epoch 156); seed 202 reached best 7.47 at step 2{,}562 (stopped at step 4{,}662, epoch 173). The step-faithful protocol actually yields \emph{lower} best dev BLEU than the epoch-validation approximation (7.0--7.5 vs 8.6--9.5), because more frequent validation accelerates the plateau-scheduler's LR decay. On PHX-public, both show near-zero gaps ($+0.02$ and $-0.30$), family-range pass-through (1.11 and 1.17), and family-range seen-holdout (5.71 and 6.59). The released evaluator reached 10.01 by step 770 on the same protocol; the reconstruction trajectory diverges after step $\approx$600, where the original accelerates while the reconstruction plateaus.

\paragraph{Frozen-protocol confirmation.}
To bound the selective-analysis risk of these post-hoc characterizations, we froze the diagnostic protocol (gap, pass-through, holdout readout, dose--response position, perturbation pattern) and then trained two new documented-recipe seeds (1506, 1607) that had not entered any prior analysis. Both confirm the family pattern exactly: gaps $-1.03$ and $-0.98$, pass-through 1.22 and 1.29, seen-holdout 6.8 and 7.0, dev BLEU 8.1 and 7.6 within the family competence range. The confirmation is narrow --- it covers the family side of the comparison, not the released side, which has no retrainable provenance.

\subsection{Checkpoint-targeted oracle stress tests (appendix-diagnosed)}
\label{sec:non-bleu-audit}
\label{sec:bt_likelihood_attack}

We treat evaluator-aware donor selection as a checkpoint-targeted oracle stress test, not as a deployable SLP method, and we report its unmasked advantage as a non-robust diagnostic: the diagonal advantage does not survive reference-overlap exclusion (Appendix~\ref{sec:appendix_oracle}). The full cross-checkpoint transfer matrix, candidate-pool sweep, and selector/evaluator-label permutation test are likewise relegated to the appendix because they share target references and evaluator by design. The main text states only that the unmasked advantage is not robust to reference-overlap exclusion and that the white-box search budget (candidate-pool size) controls whether any decoded-BLEU margin is detectable.


\section{Discussion}
Our results characterize a released-evaluator stress-test susceptibility and, more informatively, the property that produces it. The original retrieval--reference reversal does not reproduce across fourteen recipe-conditioned reconstruction runs (seed-level mean gap $-0.55$ BLEU points, 95\% prediction interval $[-1.4,+0.3]$), and it does not reproduce across train-pool scaling (12.5\%--75\%), twenty rescue variants across six hyperparameter families, all three examined checkpoint-selection objectives (NLL, decoded BLEU, decoded WER) over the dense training trajectories, or two step-faithful runs following the corrected released protocol (\S\ref{sec:stepfaithful}). That is a broad negative result, but the more useful finding is positive: the released evaluator exhibits a measurable signature that no constructible reconstruction shows (\S\ref{sec:positive_evidence}). It shows a strong training-pool readout signature, it reproduces donor transcripts from donor motion far beyond the family's competence trend, and its replay advantage tracks that readout: reversing donor frames removes 16 BLEU points, restricting donors to unseen test poses turns the gap negative, and replacing the speech-transcribed references with human sign-to-text references attenuates the gap by about 90\% descriptively (\S\ref{sec:human_refs}). Within the observed competence range (dev BLEU-4 4.5--10.0) the gap--competence response is flat; because no reconstruction reached the released evaluator's competence (13.38), nonlinear behavior in the unsupported interval cannot be excluded. The evidence is consistent with the anomaly being a property of the released checkpoint and its data pipeline --- which, unlike every reconstruction protocol we can build from public information, produces strong training-content readout --- interacting with the original reference frame. Whether that difference reflects unpublished training history, data handling, checkpoint selection, or other undocumented details cannot be identified from the released artifact alone.

These findings extend prior critiques of BT validity from output faithfulness to ranking stability: improved BT and distributional scores need not imply target conditioning \cite{hong2026conditional} (recent preprint), low recorded-pose BT baselines are documented \cite{walsh2024datadriven}, and BT likelihood is a more stable correlate of human judgment than decoded BLEU \cite{jiang2025meaningful}. Our reversal is a \emph{decoded-sacreBLEU} reversal; teacher-forced BT likelihood does not reproduce it (GT NLL $3.060$ vs PURE $3.479$ per token), consistent with that divergence. Whole-donor replay is a constructed probe, whereas SOKE-style word-level retrieval can improve pose quality under a different representation and metric \cite{soke}; and multiple text scorers applied to the same BT decoding are not independent confirmation, because they inherit the same recognizer output.

Where feasible, learned BT studies should test sensitivity across independently trained evaluator checkpoints, combine controls and uncertainty estimates with both evaluator replication and retrieval-pool substitution, report decoded-text metrics alongside teacher-forced BT likelihood, and audit reference validity with human sign-to-text references where available \cite{czehmann2026sacred}, particularly when systems are tuned against evaluator feedback or draw donors from the evaluator's training pool.

\section{Limitations}
\label{sec:limitations}
The primary evidence is confined to PHX-public, a template-dense weather domain with documented reference-validity and train--test-overlap issues~\cite{czehmann2026sacred,alkain2026critical}; CSL-Daily is protocol-sensitive. Results should not be generalized to broader domains, datasets, pose representations, or evaluator architectures.

The fourteen recipe-conditioned reconstructions are less competent than the released evaluator (dev BLEU-4 differences 4.3--6.7 points; all rescue, config-faithful, and long-schedule maxima 10.0), and the original training history is unpublished, so competence-matched replication was not possible. The reconstructions are recipe-conditioned --- they share the documented architecture, tokenizer, manifests, and selection rule but are not IID repeats of the original training process --- and the fourteen-seed prediction interval describes only that population. The positive characterization of \S\ref{sec:positive_evidence} (readout signature, pass-through, flat dose--response) does not require competence parity, but it still cannot identify \emph{which} unpublished training detail produces the released evaluator's readout property.

The donor-pool substitution decomposition is path-dependent (Table~\ref{tab:path_sensitivity}); the pool-origin contrasts and the factorial main effect are descriptive associations, not identified training-exposure effects, because train-pool and test-pool donors differ in recording conditions and template composition, signer-independent generalization is poor on this corpus~\cite{artiaga2025signer}, and the factorial strata retain a similarity SMD of $+0.4$ between seen and unseen donors (Table~\ref{tab:factorial_design}). The reference-replacement attenuation ($\approx$90\%) is descriptive, not a causal attribution; the human back-translations are a single-annotator derivative resource (one deaf fluent L2 DGS signer) whose high-confidence subset lacks second-annotator adjudication.

The reversal is a \emph{decoded-BT-text-metric} reversal, not a sign-language quality reversal: the decoded-text metrics agree because they consume the same beam-decoded hypothesis, teacher-forced BT likelihood does not reproduce it, and the transcript-slot diagnostic (\S\ref{sec:slot_f1}) is reference-anchored and rule-based --- it observes no spatial grammar or non-manual content, and its extractor has moderate recall (0.86 on our 50-item audit, single auditor). Without target-conditioned DGS signer judgment or a validated sign-native pose metric, we cannot determine whether the retrieved donor motion conveys target-equivalent meaning; we explicitly do not claim that donor motion fails to express the target, nor that the reversal indicates dataset contamination or pose memorization. The SiLVERScore implementation accepts RGB video or I3D features, not pose tensors, and the SignCLIP-style 203-landmark input cannot be unambiguously mapped from the released 178-point skeleton, so we do not run a sign-native metric under an unvalidated mapping (unvalidated comparability, not intrinsic incompatibility \cite{silverscore,jiang2025meaningful}).

For completeness, we also computed BLEURT scores with \texttt{Elron/bleurt-large-512} (a PyTorch-native substitute for BLEURT-20, whose TF checkpoint fails under our setup): they agree in direction with the other decoded-text metrics (gap $+0.165$ under the original, near zero or negative under reconstructions). Rank correlation with canonical BLEURT-20 is high ($r=0.87$, $\rho=0.81$ on 100 pre-registered German pairs) but absolute scores are biased ($-0.47$), so we report them in the artifact only, as rank-based sensitivity, and not as comparable absolute values. A limitation we acknowledge openly: no Deaf community member was involved in the audit design, and no DGS-based semantic judgment was collected; Deaf-led reviews identify this absence as a structural risk in sign-language AI~\cite{slai_biases}.

\section{Conclusion}
The released SLRTP2025 back-translation evaluator scores a constructed whole-donor replay stress test (TN-PURE-v1, not an actual SLRTP2025 leaderboard entry) at 23.79 sacreBLEU against 12.78 for the recorded test poses on PHX-public, a $+11.01$ BLEU-point decoded-sacreBLEU reversal invariant to the scorer version. This audit establishes three findings. First, the reversal does not reproduce: fourteen recipe-conditioned reconstructions, train-pool-scaling variants, twenty rescue runs, all three examined checkpoint-selection objectives (NLL, decoded BLEU, decoded WER) over the training trajectories, and two step-faithful runs under the corrected released protocol all yield gaps within $[-1.1,+0.3]$ BLEU points. Second, it is highly reference-dependent: it attenuates from $+11.01$ to $+1.12$ BLEU points against public human sign-to-text back-translations (descriptive attenuation, not a causal attribution). Third, and most informatively, the released evaluator is positively characterized by a readout-with-generalization signature that no constructible reconstruction exhibits: it decodes its full training pool at 78.4 BLEU with a 69.4\% exact-match rate while generalizing to unseen data (13.38/12.78), reproduces donor transcripts at 2.4$\times$ the family pass-through trend (3.09 vs 1.05--1.31), loses $-16.1$ points under frame reversal, and its advantage grows with donor--query similarity. The reversal is confined to decoded BT-text metrics; teacher-forced BT likelihood does not reproduce it. The evidence is consistent with the anomaly being a property of the released checkpoint and its data pipeline, co-occurring with unusually strong train-pool readout under the original reference frame; the specific mechanism is not identifiable. It is not an evaluator-identity claim, not a leaderboard-rank statement, and not evidence that a randomly reseeded reconstruction would behave alike. Learned BT studies should combine evaluator replication, retrieval-pool substitution, reference-validity audits with human sign-to-text references, and teacher-forced likelihood reporting, particularly when systems are tuned against evaluator feedback or draw donors from the evaluator's training pool.
\appendix

\section{Analysis Provenance and Multiplicity}
\label{sec:provenance}
Table~\ref{tab:provenance} classifies every analysis in this manuscript by when and why it entered the paper, its evidentiary status, and its multiplicity handling. The single primary estimand is the PURE--GT decoded-sacreBLEU gap under the released evaluator (\S\ref{sec:method}); everything else is secondary, and secondary results are either Holm-corrected within a defined family or explicitly descriptive.

\begin{table}[h]
\centering
\scriptsize
\caption{Analysis provenance. ``Prespecified'' = fixed before the headline result was observed; ``Post-hoc'' = added during the audit; ``Correction'' = supersedes an earlier pipeline error. C = confirmatory (frozen protocol, error control), D = descriptive (no formal error control).}
\label{tab:provenance}
\begin{tabular}{p{6.2cm}p{2.8cm}p{0.9cm}p{3.0cm}}
\toprule
Analysis & Provenance & Status & Multiplicity handling \\
\midrule
PURE--GT gap, original evaluator (primary estimand) & Exploratory-observed; protocol frozen afterwards & D & descriptive MC estimate $p\leq10^{-4}$ \\
Protocol checks (scorer, FPS, decoding sweep) & Prespecified & C/D & exact-match / descriptive \\
Conditioning controls (retrievers, gloss, random) & Prespecified & D & descriptive \\
14 reconstruction gaps + prediction interval & 6 prespecified; 8 post-hoc & C/D & DiD CIs; PI descriptive \\
Cross-metric gap table & Correction (metric alignment) & D & direction/sign only \\
Template-density strata; donor-exclusion $\tau$-sweeps & Prespecified & D & descriptive \\
Template-holdout, BT-retrain, cross-fit & Prespecified & C & paired bootstrap CIs \\
Donor-pool substitution; canonical B/C/D path & Post-hoc & D & query-index bootstrap CIs \\
Alternative paths, Shapley, S4 system; balance/caliper/signer & Post-hoc & D & path-sensitivity display \\
605-query factorial (cells, effects, interaction) & Prespecified design & C & Holm across 3 contrasts \\
Reference replacement (single/dual); overlap strata & Post-hoc & D & paired bootstrap of gap difference \\
Transcript-slot diagnostic + extractor audit & Post-hoc & D & descriptive; extractor P/R \\
Five related diagnostic views (\S\ref{sec:positive_evidence}) & Post-hoc & D & triangulation, not independent \\
Oracle/white-box transfer matrix & Post-hoc; appendix-only & D & permutation check \\
Competence gate, trajectories, Pareto, rescue (20); gate threshold sensitivity & Gate prespecified; rescue expanded & D & gate counts; 12-threshold grid \\
Big-arch variants; long-schedule runs & Post-hoc & D & competence attempts \\
In-sample readout diagnosis (train/dev/test) & Author-initiated & D & descriptive; no identity claim beyond released train.pt \\
Rank stability (9 systems $\times$ 7 evaluators) & Post-hoc & D & $\tau_b$/$\rho$/flip means \\
IPW pool-origin robustness & Post-hoc & D & balance robustness check \\
Slot hypothesis-side audit & Post-hoc & D & extractor P/R by system \\
Frozen-protocol confirmation (seeds 1506/1607) & Post-hoc & C (frozen protocol) & confirmation of family-side pattern \\
Ethics/licensing declarations & Updated each revision & --- & --- \\
\bottomrule
\end{tabular}
\end{table}

\section{Checkpoint Registry}
\label{sec:checkpoint_registry}
Table~\ref{tab:checkpoint_registry} is the definitive registry of every BT evaluator checkpoint used in this manuscript, with the analyses each family enters; all checkpoint counts in the text refer to these families.

\begin{table}[h]
\centering
\scriptsize\setlength{\tabcolsep}{3pt}
\caption{Definitive checkpoint registry: 51 BT evaluator checkpoints trained by us, plus the released original. ``Analyses'' lists where each family is used (figure/table/statistic).}
\label{tab:checkpoint_registry}
\begin{tabular}{p{4.6cm}p{2.0cm}p{7.2cm}}
\toprule
Family & Count & Enters \\
\midrule
Reconstructions (primary seeds 101--606) & 6 & multiseed table, all cp1--6 cells, bootstrap, dose--response \\
Reconstructions (extension 707--1405) & 8 & multiseed extension range, prediction interval, dose--response \\
Rescue lr (A/B $\times$ 4 seeds) & 8 & rescue table, gate counts \\
Rescue expanded (bs/dropout/wd0/ls/ep600) & 12 & rescue table; wd0 also dose--response, perturbation, pass-through, in-sample diagnosis \\
Train-pool ladder (12.5--75\%) & 4 & dose--response figure, gap ranges \\
Config-faithful (3000-epoch BLEU-selected) & 4 & config-faithful table, dose--response \\
Large-architecture variants (A1--A4) & 4 & competence-matched attempt (\S\ref{sec:config_faithful_results}) \\
Confirmation seeds (1506, 1607) & 2 & frozen-protocol confirmation \\
Long-schedule (800-epoch, BLEU-selected) & 2 & competence ceiling (dev 10.02) \\
BT-retrained (train-core only) & 1 & holdout series \\
Cross-fit BT\_A/B & 2 & holdout series \\
\midrule
Released original & 1 & audit target (not trained by us) \\
\bottomrule
\end{tabular}
\end{table}

\section{Supplementary Oracle Diagnostics}
\label{sec:appendix_oracle}
\paragraph{Source-conditioned retrieval is legal, oracle-gloss retrieval is leakage.}
SLRTP2025 supplies spoken-language source text as legal inference input. Under pure whole-donor copy, source-text Jaccard scores 23.8 sacreBLEU (0--100 scale, as throughout this appendix), LaBSE 19.2, and multi-SBERT 16.5; the PT-composed source-text condition scores 18.9 (Table~\ref{tab:isomorphic_conditioning}). Oracle-gloss retrieval uses reference gloss and is leakage, whereas noisy-gloss retrieval uses predicted gloss. Thus the strongest legal source-conditioned result is the 23.8 whole-donor stress test, which uses no generator or scaffold.
\paragraph{Disclosure of selection/evaluation reuse.}
The white-box oracle analyses in this appendix reuse the same target reference and evaluator for both donor selection and decoded evaluation. This is an intentional adversarial design: diagonal entries quantify susceptibility to a selector that knows its evaluator, not an independent validation of retrieval quality.

\paragraph{Cross-checkpoint transfer matrix.}
For each query, a selector maximizes its evaluator's teacher-forced likelihood over candidate donor poses; the selected fixed pose is then decoded by every evaluator (token-normalized NLL, sorted donor IDs, \texttt{numpy.argmax} first tie-break). All seven diagonal margins are positive (mean $+0.0150$, 95\% paired-bootstrap CI $[+0.0105,+0.0198]$). Original-selected donors transfer to five of six reconstructions (mean local-GT margin $+0.0074$, $[+0.0016,+0.0135]$); the reverse direction fails in all six cells (reconstruction-selected donors evaluated by Original, mean margin $-0.0455$ $[-0.0554,-0.0358]$). The full 7$\times$7 matrix, unrounded cells, margins, and CIs are in the artifact.



\paragraph{Candidate-pool size.}
The Original-selector stress test shows a strong search-budget effect: at 8, 32, 128, 512 and 2,048 candidates, decoded sacreBLEU remains below local GT by $-12.1$, $-9.2$, $-5.7$, $-3.6$ and $-1.3$ BLEU points respectively; only the full 7,060-donor pool is detectably positive ($+1.7$, $[+0.4,+3.0]$).

\paragraph{Reference-overlap exclusion.}
Masking 1,036 exact-text query--donor pairs before selection: Original-on-Original decoded BLEU is 13.86 versus local GT 12.78 (absolute difference $+1.08$ BLEU points), but the paired relative-margin 95\% CI crosses zero; the frozen joint likelihood-and-interval success rule is met in 0/7 checkpoints. Masking all 19,694 pairs with source-token Jaccard $\geq 0.5$: the score falls to 8.30 ($-4.47$), again 0/7 meet the rule. The full-pool diagonal advantage does not survive as a statistically robust decoded-BLEU result after exact-text exclusion and disappears after high-overlap exclusion.



\paragraph{Noisy-gloss retrieval is split-dependent.}
A 6-layer text$\to$gloss Transformer (dev BLEU-4 13.17) yields noisy-gloss retrieval sacreBLEU 12.7 on PHX-public (strong-NMT), versus 1.6 for a weak NMT (BLEU-4 1.15) and 6.1 for XLM-R embedding retrieval; donor exclusion reduces the strong-NMT value to 10.9 at $\tau{=}0.5$, and the same control reaches 0.4 on the template-disjoint holdout (Table~\ref{tab:holdout_controls}).




\paragraph{PHX-public pure-donor donor-exclusion $\tau$-sweep.}
\label{sec:appendix_public_donor_exclusion}
To test whether the headline reversal survives donor exclusion on the same split, we applied a source-text Jaccard $\tau$-sweep and exact-text-exclusion to the canonical pure text-nearest donor on PHX-public under the original SLRTP2025 BT (Table~\ref{tab:pub_tau_sweep}). The reversal survives exact-text exclusion (23.0 vs GT 12.8, 32/641 donors replaced) and source-Jaccard exclusion up to $\tau{=}0.5$ (17.8 vs GT 12.8); only at $\tau{=}0.3$ and $\tau{=}0.1$ does the retrieval-vs-GT gap turn negative.

\begin{table}[t]
\centering
\small
\caption{PHX-public pure-donor source-text donor-exclusion $\tau$-sweep under the original SLRTP2025 BT (canonical tie-breaking, 12.5\,fps; sacreBLEU on the 0--100 scale). The retrieval-exceeds-GT reversal survives exact-text exclusion and $\tau{=}0.5$ exclusion.}
\label{tab:pub_tau_sweep}
\begin{tabular}{lrrr}
\toprule
Condition & sacreBLEU & $\Delta$ vs GT & Reversal? \\
\midrule
GT baseline                             & 12.8 & ---   & --- \\
Text-nearest pure (no exclusion)        & 23.8 & $+11.0$ & YES \\
Exact-text excluded                     & 23.0 & $+10.2$ & YES \\
$\tau{=}0.7$ (Jaccard $>$ 0.7 excluded) & 21.2 & $+8.4$ & YES \\
$\tau{=}0.5$ (Jaccard $>$ 0.5 excluded) & 17.8 & $+5.0$ & YES \\
$\tau{=}0.3$ (Jaccard $>$ 0.3 excluded) & 8.0 & $-4.8$ & NO \\
$\tau{=}0.1$ (Jaccard $>$ 0.1 excluded) & 0.2 & $-12.5$ & NO \\
Random donor (mean of 20 seeds)          & 0.9 & $-11.6$ & NO \\
\bottomrule
\end{tabular}
\end{table}

As a same-paradigm proxy for the SLRTP2025 winner (whose checkpoints are not public), the noisy-gloss retrieval+LM system shows the same pattern: it survives exclusion up to $\tau{=}0.5$ (10.9, still $6.8\times$ the PT-only baseline and $11.7\times$ the random-donor mean; mean selected-donor Jaccard 0.46) and collapses only at $\tau{=}0.1$.


\section{Supplementary Template-Density and Holdout Boundary Analyses}
\label{sec:appendix_template_boundary}
\paragraph{Template density does not gate the reversal (within-pipeline test).}
On the public split (same SLRTP2025 evaluator, 178-layout, full-sequence), we stratify the 641 test sequences by template density (gloss bigram-Jaccard with nearest train donor) and compute corpus sacreBLEU per stratum:

\begin{table}[h]
\centering
\small
\caption{Template-density stratification on the public split (canonical tie-breaking; sacreBLEU on the 0--100 scale). Retrieval exceeds GT at \emph{all} density levels. The high-density gap is 1.2$\times$ the low-density gap; density is a moderate moderator of magnitude, not a gate on presence.}
\label{tab:template_density_strata}
\begin{tabular}{lrrrr}
\toprule
Density stratum & Mean density & GT & Retrieval & Gap \\
\midrule
Low  (n=249)  & 0.152 & 5.9 & 16.6 & $+10.7$ \\
Mid  (n=179)  & 0.254 & 9.5 & 19.7 & $+10.1$ \\
High (n=213)  & 0.510 & 26.1 & 38.6 & $+12.4$ \\
\bottomrule
\end{tabular}
\end{table}

The reversal is present even in the lowest-density stratum ($+10.7$ BLEU points), so membership in the high-density stratum is not necessary. The high-density gap ($+12.4$) is 1.2$\times$ the low-density gap ($+10.7$). The sample-level Spearman correlation is weak-to-moderate ($\rho{=}{+}0.29$); we do not interpret $\rho^2$ as variance-explained because the relation is non-linear. The frozen holdout's mean density is 0.188 (median 0.167), between the public split's low (0.152) and mid (0.254) strata, so the holdout is not a dramatically different data regime.

\paragraph{Granularity, same-evaluator holdout, BT-retrain, and cross-fit.}
Truncating public-split sequences to 64 and 32 frames shrinks the gap from $+11.0$ (full) to $+10.5$ and $+5.6$ BLEU points, but the reversal persists at all frame caps: granularity amplifies but does not gate. On a template-holdout built from SLRTP training data (178-layout, same evaluator), GT rises from 12.8 to 77.0 because the released BT was trained on all 7{,}060 train sequences including the 525 holdout sequences (readout signature, cf. \S\ref{sec:positive_evidence}). Retraining the BT model from scratch on train-core only (6{,}535 sequences, holdout excluded) collapses holdout GT to 6.2 and retrieval to 5.4 (no reversal); applying that retrained BT to PHX-public drops canonical donor copy from 23.8 to 8.3 and GT from 12.8 to 8.7 (gap $-0.4$, not detectable). A donor-familiarity cross-fit (hash-split halves A/B, each BT trained on the other half only, never seeing its evaluation half) gives pooled GT 4.0 vs retrieval 3.9 (gap $-0.1$; with $n{=}659$ the minimum detectable effect is uncalibrated, so we read ``no large donor-familiarity effect detected''). Table~\ref{tab:holdout_series} summarizes the four controls; none reproduces the original reversal.

\begin{table}[h]
\centering
\scriptsize
\caption{Holdout boundary controls (sacreBLEU, 0--100 scale). None of the four controls reproduces the original $+11.0$ reversal; all confound training exposure, competence, and pipeline, so they bound rather than identify the checkpoint property.}
\label{tab:holdout_series}
\begin{tabular}{llrrr}
\toprule
Control & Setup & GT & Retrieval & Gap \\
\midrule
Same evaluator & test split (template-overlap) & 12.8 & 23.8 & $+11.0$ \\
Same evaluator & template-holdout (seen in training) & 77.0 & 20.1 & $-56.9$ \\
BT-retrained (unseen holdout) & train-core 6{,}535 & 6.2 & 5.4 & $-0.9$ \\
Retrained BT on PHX-public & unseen to both & 8.7 & 8.3 & $-0.4$ \\
Cross-fit A/B (donor-unseen) & pooled, $n{=}659$ & 4.0 & 3.9 & $-0.1$ \\
\bottomrule
\end{tabular}
\end{table}

Together with the fourteen-run factorial, these controls bound the reversal to the examined original checkpoint and pipelines; they do not identify its causal checkpoint property.

\section{Supplementary CSL-Daily Boundary Control}
\label{sec:csl-daily}
CSL-Daily (Chinese Sign Language, daily-life discourse) serves as a non-template-domain control, but its recognizer is too weak to support robust claims. Three distinct protocols appear in the literature and in our data, and we keep them separate: (i) the 256-window protocol (per-window sacreBLEU 0.2--0.6 with GT at 0.36 and 40.2\% empty predictions); (ii) the 47-sequence subset spanned by those windows (aggregated GT BLEU-1 70.8 versus 10.9--14.6 for retrieval); and (iii) the full 1{,}176-sequence Mode-B reconstruction (GT BLEU-1 79.1, sacreBLEU 26.7). Because the window and sequence protocols differ qualitatively, CSL-Daily remains a protocol-fragility boundary rather than negative shortcut evidence; we report it only as a boundary condition, not as external validation. Under Mode B sequence-level aggregation (Table~\ref{tab:csl_donor_exclusion}), text-nearest donor copy reaches only 6\% of GT sacreBLEU (1.6 vs 26.7), and the $\tau$-sweep shows the expected monotonic trend (donors with higher Jaccard overlap yield slightly higher BLEU-1).

\begin{table}[t]
\centering
\small
\caption{CSL-Daily Mode B donor-exclusion results (1{,}176 sequences; 0--100 scale). Text-nearest donor copy (sacreBLEU 1.6) is far below GT (26.7). CSL-Daily is treated as inconclusive (recognizer protocol sensitivity) rather than as negative evidence.}
\label{tab:csl_donor_exclusion}
\begin{tabular}{lrr}
\toprule
Condition & sacreBLEU & BLEU-1 \\
\midrule
GT baseline & 26.7 & 79.1 \\
Text-nearest donor & 1.6 & 43.5 \\
Random donor & 0.0 & 3.9 \\
$\tau{=}0.5$ exclusion & 1.0 & 42.6 \\
BM25 retrieval & 1.5 & 39.2 \\
TF-IDF retrieval & 1.2 & 38.5 \\
\bottomrule
\end{tabular}
\end{table}

\section{Supplementary Donor and Template Analyses}
\label{sec:donor_decomposition}
A critic might argue that text-nearest donors are semantically legitimate in a template-dense weather domain. The mean text Jaccard between query and donor transcript is 0.535 (median 0.500, 5.0\% exact), and the pass-through analysis of \S\ref{sec:positive_evidence} shows the released evaluator reads the donor's own transcript, not the query's intended meaning (3.09$\times$ ratio); these diagnostics do not establish either a causal lexical-overlap mechanism or target-equivalent DGS meaning.

\subsection{Template-holdout controls and donor-exclusion robustness}
\label{sec:frozen-template-holdout}
\label{sec:aggressive-donor-audit}
To address public-split reuse directly, we constructed a deterministic template-holdout with disjoint windows, source sequences, and exact gloss templates (14{,}988 train-core windows from 4{,}986 sequences; 1{,}538 holdout windows from 528 sequences; selection seed \texttt{rcp-confirmation-template-holdout-20260714-v1}; window-, sequence-root-, and exact-gloss-template-disjoint leakage guards), retrained every component on train-core only, and evaluated once. Text-nearest retrieval exceeds random donor copy ($+2.3$, CI $[+2.0,+2.7]$) yet stays far \emph{below} the GT baseline ($-10.8$, CI $[-11.8,-9.9]$): the shortcut is real over random, but retrieval does not exceed GT on the clean split. Donor-exclusion and alternative-retriever sweeps on the same split (Table~\ref{tab:holdout_controls}) show the retrieval advantage persisting at $\tau{=}0.5$ exclusion and under BM25/TF-IDF, consistent with sub-template lexical reuse (mean selected-donor gloss-Jaccard 0.40 at $\tau{=}0.5$), and collapsing to the random baseline only at $\tau{=}0.1$; a small non-monotonic step between $\tau{=}0.7$ and $\tau{=}0.5$ has CI including zero.

\begin{table}[t]
\centering
\scriptsize
\caption{Frozen template-holdout controls (1,538 windows, $\alpha{=}0.5$; sacreBLEU 0--100). Top: main systems. Bottom: donor-exclusion sweep and alternative retrievers. The retrieval advantage persists at $\tau{=}0.5$ and under BM25/TF-IDF, reaching the random baseline only at $\tau{=}0.1$.}
\label{tab:holdout_controls}
\begin{tabular}{lrrr}
\toprule
System & sacreBLEU & Empty \% & Mean Jac. \\
\midrule
GT baseline & 13.10 & 9.95 & --- \\
Random train donor & 0.04 & 11.83 & --- \\
Text-nearest train donor & 2.30 & 10.66 & --- \\
LaBSE / multi-SBERT (donor copy) & 1.73 / 1.51 & 9.2 / 11.3 & --- \\
Noisy-gloss retrieval (strong NMT) & 0.40 & 10.90 & --- \\
Retrieval composition ($\alpha{=}0.5$) & 2.20 & 10.79 & --- \\
\midrule
$\tau{=}0.1$ / 0.3 / 0.5 / 0.7 exclusion & 0.04 / 1.05 / 1.93 / 2.01 & 10.0--10.5 & 0.10/0.29/0.40/0.42 \\
BM25 / TF-IDF nearest & 1.75 / 2.36 & 10.8 / 8.8 & --- \\
\bottomrule
\end{tabular}
\end{table}

\paragraph{Template-density context.}
The SLRTP2025 test vocabulary is nearly closed: 95.4\% of gloss types and 99.6\% of gloss tokens occur in train (German text: 94.1\% and 99.3\%). Nearest donors share mean gloss unigram containment 0.557 and bigram-Jaccard 0.302 (German text: 0.610 and 0.322). Progressive Transformer outputs show Spearman $-0.06$ ($p{=}0.12$) between donor proximity and BLEU, versus $+0.29$ ($p{<}0.001$) for retrieval; because PT sacreBLEU is 1.6, near the metric floor, this contrast is suggestive rather than definitive.

\section*{Declarations}

\paragraph{Funding.} Anonymous.

\paragraph{Conflicts of interest.} The authors declare no conflicts of interest.

\paragraph{Ethics statement.}
No new sign-language participants were recruited for this work. All experiments use publicly released datasets (PHOENIX-2014T, CSL-Daily, SLRTP2025), publicly released human sign-to-text back-translations and confidence annotations of Czehmann et al.~(2026), and rule-based text analysis applied to decoded hypotheses. The single human auditor of the slot extractor (50 references + 50 hypotheses) was one of the authors; no personally identifiable information is collected or reported.

\paragraph{Data licensing and availability.} PHOENIX-2014T (RWTH-PHOENIX-Weather 2014T) is distributed by RWTH Aachen under the Creative Commons BY-NC-SA \textbf{3.0} license shown on the official corpus page (checked 2026-06-10); we comply with its non-commercial/share-alike terms and do not redistribute the corpus. The human sign-to-text back-translations and annotations of Czehmann et al.~\cite{czehmann2026sacred} are a \emph{derivative} resource distributed separately under CC BY-NC-SA \textbf{4.0}, which covers only their annotations/back-translations, not the original corpus; we do not redistribute either. CSL-Daily and the SLRTP2025 pose bundle are likewise obtained from their providers under provider-specific terms. We did not find local documentation establishing permission to redistribute signer-derived pose tensors, rendered signer motion, or videos. The anonymous artifact therefore contains sequence IDs, exclusion manifests, hashes, code, evaluator configurations, and derived sufficient statistics, but not source pose/video data. Pose tensors may retain signer identity and style information even after skeletonization; we release no per-signer raw poses. Users reconstructing the audit must obtain PHOENIX-2014T, CSL-Daily, and SLRTP2025 inputs through their official access channels.

\paragraph{Code availability.} The anonymous review artifact is available at \url{https://anonymous.4open.science/r/rcp-audit-artifact-B314/README.md} (individual files are accessible; the landing page may return an HTTP 401 for some browsers --- use the README.md link directly). From its root, \texttt{make core-audit PYTHON=/path/to/python3} reruns checkpoint--GT alignment regression checks and the 60-cell BLEU decomposition. It contains all 38 analysis scripts, 153 result files, decoded hypotheses for all main-text cells, per-query donor registries, the machine-readable checkpoint registry (51 checkpoints with dev metrics and gate status), SHA-256 file digests, dose-response and trajectory figures, and sufficient statistics for every bootstrap and permutation test. A text-only reproduction path (no pose data) can regenerate every text-metric number and confidence interval from the decoded cells and result files.

\paragraph{Computational cost.} Experiments ran on one node with eight NVIDIA GeForce RTX 2080 Ti GPUs. Per-family cost (GPU-hours, from epoch logs and run timestamps): the six primary evaluator reconstructions 1.4 (0.20--0.26 each); the eight extension seeds 1.81 (0.204--0.262 each); the twenty competence-rescue runs approximately 5.0 (0.19--0.31 each); fixed-output decoding of the canonical four systems approximately 2 minutes per evaluator; the 605-query factorial re-scoring 0.17 total across seven evaluators; donor-registry construction, matching variants, and bootstrap analyses are CPU-bound (approximately 40 CPU-hours for the 10{,}000-resample decompositions). Training the KL-VAE requires approximately 6 minutes (100 epochs, single GPU) and the text-conditioned diffusion generator approximately 1 hour (100 epochs, single GPU) on the template-holdout split. Back-translation evaluation per system row takes 5--20 minutes depending on recognizer and window count. The two-week wall-clock envelope also includes sequential ablations, data preparation, failed protocol checks, and idle queue time, so it is not reported as GPU-hours.

\bibliographystyle{sn-basic}
\bibliography{references}

\end{document}
